% !TeX program = xelatex
\documentclass[11pt,a4paper,fontset=fandol]{ctexart}
\usepackage[margin=18mm]{geometry}
\usepackage{amsmath,amssymb,bm,enumitem,hyperref}
\hypersetup{hidelinks,pdftitle={大一数学分析期末复习·上册}}
\setlist[enumerate]{label=（\arabic*）,leftmargin=2.6em,itemsep=0.25em}
\newcommand{\dd}{\,\mathrm{d}}
\newenvironment{exercise}[2]{%
  \par\medskip\noindent\textbf{第#1题}\quad\textit{#2}\par\smallskip
}{\par\medskip}
\allowdisplaybreaks
\setcounter{secnumdepth}{-1}
\sloppy
\setlength{\emergencystretch}{4em}
\title{大一数学分析期末复习·上册\\[0.4em]\large 基础计算与常见题型}
\author{}
\date{2026年10月20日}
\begin{document}
\maketitle
\section*{使用说明}
本册用于大一期末分章复习，覆盖函数与极限、导数与微分、中值定理与导数应用、不定积分、定积分及其应用、微分方程。题目以基础计算、常见应用和课堂短证明为主，证明题使用已学定理并说明关键条件。

全册共830题，按阅读顺序统一编号1至830。40节各保留两道补充练习，用于基础变式与常规小综合，均纳入连续编号。本版设“基础”“常规”两档：“基础”表示直接运用概念、公式或标准方法，“常规”表示期末常见的两步计算、应用或短证明。建议先做基础题，再做常规题。

常规版与iPad横屏版共用同一份题面，后者按每页一道题设计，下方纯空白用于草稿。本次仅交付LaTeX源文件，按要求未编译，未附答案。未获得本校期末试卷，难度按通用的大一期末复习要求调整；具体范围以课堂教学为准，微分方程按课程实际教学选用。
\tableofcontents
\clearpage
\section{第1章 函数与极限}
\subsection{函数}
\begin{exercise}{1}{基础·定义域}
求 $f(x)=\arccos(2x+1)+\log_2(3x+2)$ 的定义域。
\end{exercise}
\begin{exercise}{2}{基础·复合函数}
设 $f(x)=\frac{x-1}{x+1}$。求 $f(-2)$ 和 $f(f(x))$，写出复合函数的定义域。
\end{exercise}
\begin{exercise}{3}{基础·换元求函数}
已知对一切 $x\ne0$，$f(1/x)=4/x+2x^2+1$。求 $f(t)$（$t\ne0$）及 $f(x^2+1)$。
\end{exercise}
\begin{exercise}{4}{基础·反三角函数}
设 $f(x)=\arcsin(4/x)$。求定义域及 $f(-4)$、$f(4\sqrt2)$。
\end{exercise}
\begin{exercise}{5}{常规·函数值域}
求 $f(x)=\sin\frac{\pi x}{2(1+x^2)}$ 的值域。提示：先求 $x/(1+x^2)$ 的取值范围。
\end{exercise}
\begin{exercise}{6}{基础·反三角函数值}
计算 $\arccos(\sqrt3/2)$、$\arctan(-1)$、$\operatorname{arccot}1$，其中 $\operatorname{arccot}$ 取值于 $(0,\pi)$。
\end{exercise}
\begin{exercise}{7}{常规·换元求函数}
已知 $f(x+1/x)=\frac{4x^2}{x^4+1}$（$x>0$）。令 $t=x+1/x$，求 $t$ 的取值范围及相应的 $f(t)$。
\end{exercise}
\begin{exercise}{8}{基础·分段函数}
设 $f(x)=\begin{cases}\frac{x}{\pi}|\sin x|,&|x|<\pi/3,\\0,&|x|\ge\pi/3.\end{cases}$ 求 $f(\pi/6)$、$f(-\pi/6)$、$f(2)$，判断其奇偶性。
\end{exercise}
\begin{exercise}{9}{常规·反解复合函数}
设 $\tan(2g(x))=x^2-1$ 且 $|g(x)|<\pi/4$。求 $g(x)$ 的表达式。
\end{exercise}
\begin{exercise}{10}{基础·分段表达式}
将 $\varphi(x)=\max\{x,x^2,1\}$ 写成分段函数，并求 $\varphi(-2)$、$\varphi(0)$。
\end{exercise}
\begin{exercise}{11}{基础·复合次序}
设 $f(x)=\begin{cases}-1,&x<0,\\1,&x\ge0,\end{cases}$、$g(x)=2x-1$。分别写出 $f(g(x))$、$g(f(x))$ 的分段表达式。
\end{exercise}
\begin{exercise}{12}{常规·分段复合函数}
设 $f(x)=2+|x|$、$g(x)=\begin{cases}x,&x<0,\\x^2,&x\ge0.\end{cases}$ 求 $f(g(x))$ 和 $g(f(x))$。
\end{exercise}
\begin{exercise}{13}{基础·分段函数运算}
设 $f(x)=\begin{cases}2x+1,&x<0,\\x^2+1,&x\ge0,\end{cases}$、$g(x)=x+1$。写出 $f+g$ 与 $fg$ 的分段表达式。
\end{exercise}
\begin{exercise}{14}{常规·反函数}
设 $f(x)=\frac{1-\sqrt{1+x}}{1+\sqrt{1+x}}$（$x\ge-1$）。求其反函数及反函数的定义域。
\end{exercise}
\begin{exercise}{15}{基础·奇偶性}
判断 $(x^2+1)\sin^3x$、$\ln(x+\sqrt{x^2+1})$、$\frac{2^{-x}-2^x}{2^{-x}+2^x}$ 的奇偶性。
\end{exercise}
\begin{exercise}{16}{常规·复合函数奇偶性}
设 $f,g:\mathbb R\to\mathbb R$，$f$ 为奇函数，$g$ 为偶函数。由定义判断 $f\circ f$、$f\circ g$ 的奇偶性。
\end{exercise}
\begin{exercise}{17}{常规·单调性的定义}
不用导数证明 $f(x)=x^3+2x+3$ 在 $\mathbb R$ 上严格递增。提示：对 $x_1<x_2$ 考察 $f(x_2)-f(x_1)$。
\end{exercise}
\begin{exercise}{18}{基础·有界性}
判断 $f(x)=1/(x-2)$ 在 $(2,3)$ 和 $(3,+\infty)$ 上是否有界。
\end{exercise}
\begin{exercise}{19}{基础·周期}
求 $f(x)=2+\cos^2(2x-1)$ 的最小正周期。
\end{exercise}
\begin{exercise}{20}{常规·取整函数}
记 $\lfloor x\rfloor$ 为不超过 $x$ 的最大整数。求 $f(x)=x-\lfloor x\rfloor$ 的值域及 $f(1.2)$、$f(-1.2)$，验证 $1$ 是它的周期。
\end{exercise}
\begin{exercise}{21}{基础·函数关系}
设 $f:(0,+\infty)\to\mathbb R$ 满足 $f(x)+2f(1/x)=x$。求 $f(x)$。提示：将 $x$ 换成 $1/x$ 后联立。
\end{exercise}
\begin{exercise}{22}{常规·函数综合}
设 $f(x)=(x-1)/(x+1)$。求定义域、零点和反函数，并写出反函数的定义域。
\end{exercise}
\subsection{数列的极限}
\begin{exercise}{23}{常规·极限定义}
用 $\varepsilon$--$N$ 定义证明 $\lim_{n\to\infty}n/(2n+1)=1/2$。提示：计算它与 $1/2$ 的差的绝对值。
\end{exercise}
\begin{exercise}{24}{基础·幂数列}
判断数列 $(1/2)^n$、$(-1/2)^n$、$(-1)^n$ 是否收敛，收敛时求极限。
\end{exercise}
\begin{exercise}{25}{基础·幂数列极限}
分别求 $\lim_{n\to\infty}n^{-2}$、$\lim_{n\to\infty}1/\sqrt n$、$\lim_{n\to\infty}n/(n+1)$。
\end{exercise}
\begin{exercise}{26}{基础·奇偶项极限}
设 $x_{2k-1}=-1/(2k)$、$x_{2k}=2^{-k}$（$k\ge1$）。分别求奇、偶子列的极限，并判断 $x_n$ 是否收敛。
\end{exercise}
\begin{exercise}{27}{基础·分段数列}
设 $x_n=(n+1)/(n+2)$（$n$ 为偶数），$x_n=1$（$n$ 为奇数）。求 $\lim_{n\to\infty}x_n$。
\end{exercise}
\begin{exercise}{28}{常规·子列判别}
设 $x_n\to-1$、$y_n\to1$，且 $z_{2n-1}=x_n$、$z_{2n}=y_n$。利用收敛数列的子列性质判断 $z_n$ 是否收敛。
\end{exercise}
\begin{exercise}{29}{基础·极限运算}
已知 $x_n\to2$。求 $|x_n|$、$x_{n+1}$、$1/x_n$ 的极限，其中倒数只考虑充分大的 $n$。
\end{exercise}
\begin{exercise}{30}{常规·振荡数列}
利用奇、偶子列判断 $x_n=(1+1/n)(-1)^n$ 是否收敛。
\end{exercise}
\begin{exercise}{31}{基础·绝对值极限}
设 $x_n\ge0$ 且 $x_n^2\to1$。求 $\lim_{n\to\infty}x_n$。提示：使用 $|x_n-1|=|x_n^2-1|/(x_n+1)$。
\end{exercise}
\begin{exercise}{32}{常规·数列极限综合}
设 $x_n=(2n+(-1)^n)/(n+1)$。求 $\lim_{n\to\infty}x_n$ 及 $\lim_{n\to\infty}(x_{n+1}-x_n)$。
\end{exercise}
\subsection{函数的极限}
\begin{exercise}{33}{常规·极限定义}
用 $\varepsilon$--$\delta$ 定义证明 $\lim_{x\to1/2}\frac{1-4x^2}{2x-1}=-2$。提示：先在 $x\ne1/2$ 时约分。
\end{exercise}
\begin{exercise}{34}{基础·无穷远极限}
分别求 $e^{2x-1}$ 与 $\arctan x$ 在 $x\to-\infty$、$x\to+\infty$ 时的极限。
\end{exercise}
\begin{exercise}{35}{基础·左右极限}
设 $f(x)=\begin{cases}(x+1)\sin x,&x<0,\\1-\cos x,&x\ge0.\end{cases}$ 求 $x=0$ 处的左右极限及双侧极限。
\end{exercise}
\begin{exercise}{36}{基础·分段极限}
设 $f(x)=\begin{cases}\ln(1-x),&x<1,\\1/\sqrt{2x+1},&x\ge1.\end{cases}$ 求 $x\to1^-$、$x\to1^+$、$x\to0$ 时的极限。
\end{exercise}
\begin{exercise}{37}{基础·分段函数极限}
设 $f(x)=\begin{cases}x+1,&x<0,\\3,&x=0,\\2-x,&x>0.\end{cases}$ 求 $f(0)$ 及 $x\to0$ 的左右极限，判断双侧极限是否存在。
\end{exercise}
\begin{exercise}{38}{常规·局部保号性}
设 $\lim_{x\to1}f(x)/x=-1$。先求 $\lim_{x\to1}f(x)$，再利用极限保号性说明 $f(x)$ 在 $1$ 的某去心邻域内为负。
\end{exercise}
\begin{exercise}{39}{常规·极限不等式}
设 $\lim_{x\to x_0}f(x)=A$、$\lim_{x\to x_0}g(x)=B$，且 $A<B$。证明在 $x_0$ 的某去心邻域内 $f(x)<g(x)$。提示：对 $g-f$ 使用保号性。
\end{exercise}
\begin{exercise}{40}{基础·振荡与夹逼}
求 $\lim_{x\to+\infty}\cos x/x$，并判断 $\lim_{x\to+\infty}\cos x$ 是否存在。
\end{exercise}
\begin{exercise}{41}{基础·局部有界性}
已知 $\lim_{x\to0}f(x)=2$。在极限定义中取 $\varepsilon=1$，证明 $f$ 在 $0$ 的某去心邻域内有界。
\end{exercise}
\begin{exercise}{42}{基础·复合极限}
已知 $\lim_{u\to0}f(u)=2$。求 $\lim_{x\to0}f(3x)$ 和 $\lim_{x\to0}f(x^2)$，其中只考虑 $x\ne0$。
\end{exercise}
\begin{exercise}{43}{常规·函数极限综合}
设 $f(x)=x\sin(1/x)+2x$（$x\ne0$）。用夹逼定理求 $\lim_{x\to0}f(x)$，并求 $\lim_{x\to0}xf(x)$。
\end{exercise}
\subsection{无穷小与无穷大}
\begin{exercise}{44}{基础·有界因子与无穷小}
用夹逼定理证明 $\lim_{x\to0}2x\sin(1/x)=0$。
\end{exercise}
\begin{exercise}{45}{基础·无穷大判定}
判断 $x^2$、$1/x$、$\sin x$ 在 $x\to+\infty$ 时是否为无穷大或无穷小。
\end{exercise}
\begin{exercise}{46}{基础·单侧极限}
分别求 $e^{1/x}$ 在 $x\to0^-$、$x\to0^+$ 时的极限，并指出在哪一过程中它是无穷小。
\end{exercise}
\begin{exercise}{47}{基础·无穷小与无穷大}
设 $f(x)=2x/(x-2)$。分别讨论 $x\to0$、$x\to2^-$、$x\to2^+$ 时的变化。
\end{exercise}
\begin{exercise}{48}{常规·乘积极限}
已知 $x_ny_n\to0$ 且 $x_n\to2$。求 $\lim_{n\to\infty}y_n$，写出使用商的极限法则的依据。
\end{exercise}
\begin{exercise}{49}{基础·有界因子估计}
求 $\lim_{x\to0}\sin x(2+\cos(3/x))$ 和 $\lim_{x\to+\infty}\arctan x/x$。
\end{exercise}
\begin{exercise}{50}{基础·无穷小判定}
判断 $1/x$（$x\to+\infty$）、$x\sin(1/x)$（$x\to0$）、$\ln(1+x)/x$（$x\to0$）是否为无穷小。
\end{exercise}
\begin{exercise}{51}{常规·极限保序性}
设 $x_n\to0$、$y_n\to1$、$z_n\to+\infty$。说明充分大的 $n$ 满足 $x_n<y_n<z_n$。提示：分别与 $1/2$ 和 $2$ 比较。
\end{exercise}
\begin{exercise}{52}{基础·由组合极限求极限}
已知 $\lim_{x\to0}(1+2x+f(x)/x)=4$。求 $\lim_{x\to0}f(x)/x$ 及 $\lim_{x\to0}f(x)$。
\end{exercise}
\begin{exercise}{53}{基础·无穷小运算}
求 $\lim_{x\to0}x^2\cos(1/x)$ 和 $\lim_{x\to0}x^2(1+\sin(1/x))$。
\end{exercise}
\begin{exercise}{54}{常规·倒数关系}
求 $\lim_{x\to0}1/x^2$ 和 $\lim_{x\to+\infty}1/(x^2+1)$，说明两式中无穷大与倒数无穷小的关系。
\end{exercise}
\subsection{极限的运算法则}
\begin{exercise}{55}{基础·极限四则运算}
已知 $f(x)\to2$、$g(x)\to3$（$x\to0$）。求 $f+g$、$fg$、$f/g$ 的极限。
\end{exercise}
\begin{exercise}{56}{基础·有限乘积极限}
求 $\lim_{x\to+\infty}(1-1/x)(2+2/x)(3-3/x)$。
\end{exercise}
\begin{exercise}{57}{基础·直接代入}
求 $\lim_{x\to1/2}(2\sin x+2\pi/x)/\cos2x$。
\end{exercise}
\begin{exercise}{58}{基础·有理化}
求 $\lim_{x\to1}(1/(x+1)-3/(x^3+1))$ 与 $\lim_{n\to\infty}(\sqrt{n+2\sqrt n}-\sqrt{n-\sqrt n})$。
\end{exercise}
\begin{exercise}{59}{基础·消因子与有理化}
求 $\lim_{x\to2}\frac{x^2+x-6}{x^2-x-2}$ 与 $\lim_{x\to0}\frac{x^2+x}{\sqrt{1+x}-1}$。
\end{exercise}
\begin{exercise}{60}{基础·最高次幂法}
求 $\lim_{x\to+\infty}\frac{2x^2+x\cos x}{x^3-x\sin x+1}$。
\end{exercise}
\begin{exercise}{61}{基础·和式极限}
求 $\lim_{n\to\infty}\sum_{k=1}^n\frac{k}{n^2+1}$ 与 $\lim_{n\to\infty}\sum_{k=1}^n2^{-k}$。
\end{exercise}
\begin{exercise}{62}{常规·望远镜求和}
求 $\lim_{n\to\infty}\sum_{k=1}^n\ln\frac{k+1}{k+2}$。提示：先合并对数。
\end{exercise}
\begin{exercise}{63}{常规·乘积对数化}
求 $\lim_{n\to\infty}\frac1{n^2}\ln\prod_{k=1}^n2^k$。
\end{exercise}
\begin{exercise}{64}{常规·有限乘积}
利用 $(1-a)\prod_{k=0}^n(1+a^{2^k})=1-a^{2^{n+1}}$，求 $\lim_{n\to\infty}\prod_{k=0}^n(1+2^{-2^k})$。
\end{exercise}
\begin{exercise}{65}{常规·单侧极限}
分别求 $\frac{x^2-1}{x+1}e^{1/(x+1)}$ 在 $x\to-1^-$、$x\to-1^+$ 时的极限。
\end{exercise}
\begin{exercise}{66}{常规·无穷远根式差}
分别求 $\sqrt{x^2+x+3}-\sqrt{x^2-x+3}$ 在 $x\to+\infty$、$x\to-\infty$ 时的极限。提示：有理化时使用 $\sqrt{x^2}=|x|$。
\end{exercise}
\begin{exercise}{67}{基础·幂指函数连续性}
求 $\lim_{x\to e}(1+\sin2x)^{\ln\ln x}$。
\end{exercise}
\begin{exercise}{68}{基础·幂指极限}
求 $\lim_{n\to\infty}(1/2)^n$ 和 $\lim_{n\to\infty}(3/2)^n$。
\end{exercise}
\begin{exercise}{69}{常规·极限反求参数}
已知 $\lim_{x\to2}\frac{x^2+ax+b}{x^2-4}=3$，求 $a,b$。提示：先令分子在 $2$ 处为零，再约分。
\end{exercise}
\begin{exercise}{70}{常规·重要极限反求参数}
求常数 $b$，使 $\lim_{x\to0}\frac{\sin2x}{\ln(1+x)}(\cos x-b)=4$。
\end{exercise}
\begin{exercise}{71}{常规·根式极限}
求 $\lim_{x\to-\infty}(4x+\sqrt{16x^2-8x-1})$。提示：有理化。
\end{exercise}
\begin{exercise}{72}{常规·根式差极限}
求 $\lim_{x\to+\infty}(\sqrt{x^4+2x^2+3}-x^2)$。
\end{exercise}
\begin{exercise}{73}{基础·极限运算变式}
已知 $f(x)\to1$、$f(x)+g(x)\to3$（$x\to0$）。求 $g(x)$ 和 $f(x)g(x)$ 的极限。
\end{exercise}
\begin{exercise}{74}{常规·根式差综合}
求 $\lim_{x\to+\infty}(\sqrt{x^2+4x+1}-x)$，并求该式在 $x\to-\infty$ 时的极限。
\end{exercise}
\subsection{极限存在准则与两个重要极限}
\begin{exercise}{75}{常规·夹逼定理}
利用 $0\le1/n!\le1/n$（$n\ge1$），求 $\lim_{n\to\infty}\sin(1/n)/n!$。
\end{exercise}
\begin{exercise}{76}{常规·和式夹逼}
求 $\lim_{n\to\infty}\sum_{k=1}^n\frac{k}{n^2+2n+k}$。提示：将每个分母分别与 $n^2+2n+1$、$n^2+3n$ 比较。
\end{exercise}
\begin{exercise}{77}{常规·振荡差式夹逼}
求 $\lim_{x\to+\infty}(\cos\sqrt{x+1}-\cos\sqrt x)$。提示：使用 $|\cos u-\cos v|\le|u-v|$，再有理化。
\end{exercise}
\begin{exercise}{78}{常规·递推数列}
设 $x_1=3$、$x_{n+1}=(x_n+4/x_n)/2$。证明 $x_n\ge2$ 且单调递减，进而求极限。
\end{exercise}
\begin{exercise}{79}{常规·单调有界数列}
设 $x_1=1$、$x_{n+1}=\sqrt{2+x_n}$。证明 $1\le x_n<2$ 且单调递增，求其极限。
\end{exercise}
\begin{exercise}{80}{常规·数列夹逼}
已知 $a_n=\int_0^1x^n\sqrt{1-x^2}\,\dd x$。利用 $0\le\sqrt{1-x^2}\le1$ 证明 $a_n\to0$。
\end{exercise}
\begin{exercise}{81}{基础·重要极限}
分别求 $\sin2x/(5x)$ 在 $x\to0$ 与 $x\to+\infty$ 时的极限。
\end{exercise}
\begin{exercise}{82}{基础·三角极限}
求 $\lim_{x\to0}\tan3x/(4x)$ 与 $\lim_{x\to0}\sin5x/(2x)$。
\end{exercise}
\begin{exercise}{83}{基础·反三角极限}
求 $\lim_{x\to0}\arcsin2x/x$ 和 $\lim_{x\to0}\arctan x/(2x)$。
\end{exercise}
\begin{exercise}{84}{基础·三角和差极限}
分别求 $\frac{2x-\sin3x}{2x+\sin4x}$ 在 $x\to0$ 与 $x\to+\infty$ 时的极限。
\end{exercise}
\begin{exercise}{85}{基础·换元极限}
求 $\lim_{x\to3}\frac{\sin(x^2-2x-3)}{x-3}$。提示：将正弦中的多项式因式分解。
\end{exercise}
\begin{exercise}{86}{基础·余弦极限}
求 $\lim_{x\to0}(1-\cos2x)/x^2$。
\end{exercise}
\begin{exercise}{87}{基础·数列重要极限}
求 $\lim_{n\to\infty}n^2\sin(3/n^2)$。
\end{exercise}
\begin{exercise}{88}{基础·第二重要极限}
求 $\lim_{x\to+\infty}(1-2/x)^{2x+1}$。
\end{exercise}
\begin{exercise}{89}{基础·商底数幂指极限}
求 $\lim_{x\to+\infty}\left(\frac{x+1}{x-1}\right)^{2x}$。
\end{exercise}
\begin{exercise}{90}{常规·二次尺度幂指极限}
求 $\lim_{x\to+\infty}\left(\frac{2x^2-1}{2x^2+1}\right)^{x^2}$。提示：把底数写成 $1$ 加无穷小。
\end{exercise}
\begin{exercise}{91}{基础·对数等价量}
求 $\lim_{x\to0}\sin3x/\ln(1+2x)$。
\end{exercise}
\begin{exercise}{92}{常规·取对数求极限}
求 $\lim_{x\to0}(1+\tan2x)^{-3/x}$。提示：先求其对数的极限。
\end{exercise}
\begin{exercise}{93}{常规·幂指极限}
求 $\lim_{x\to0}(\cos x)^{1/x^2}$。可使用 $\ln(1+u)\sim u$ 和 $1-\cos x\sim x^2/2$。
\end{exercise}
\begin{exercise}{94}{常规·指数和的极限}
求 $\lim_{x\to0}\left((2^x+8^x)/2\right)^{1/x}$。提示：取对数后使用 $a^x-1\sim x\ln a$。
\end{exercise}
\begin{exercise}{95}{基础·对数与重要极限}
求 $\lim_{n\to\infty}\ln[(1+1/n)^n]$ 和 $\lim_{n\to\infty}(1+1/n)^n$。
\end{exercise}
\begin{exercise}{96}{常规·幂指极限反求参数}
求实数 $a$，使 $\lim_{x\to+\infty}(1+a/x)^x=e^2$。
\end{exercise}
\begin{exercise}{97}{基础·夹逼变式}
求 $\lim_{n\to\infty}(2+\sin n)/n$。
\end{exercise}
\begin{exercise}{98}{常规·重要极限综合}
求 $\lim_{x\to0}(1+2x+3x^2)^{1/x}$，并说明充分小的 $x$ 下底数为正。
\end{exercise}
\subsection{无穷小的比较}
\begin{exercise}{99}{基础·等价无穷小}
当 $x\to0^+$ 时，写出 $e^{\sqrt x}-1$、$\ln(1+\sqrt x)$、$1-\cos\sqrt x$ 的等价无穷小。
\end{exercise}
\begin{exercise}{100}{基础·无穷小阶数}
当 $x\to0$ 时，写出 $\arctan2x$、$\sqrt{1+x}-1$、$1-\cos2x$ 的等价无穷小，并比较阶数。
\end{exercise}
\begin{exercise}{101}{基础·无穷远等价量}
求 $\lim_{x\to+\infty}x\frac{2x-3}{x^2-x+1}$，据此写出原分式的等价式。
\end{exercise}
\begin{exercise}{102}{基础·指数等价量}
求 $\lim_{x\to0}(e^{x^2}-1)/x^2$。
\end{exercise}
\begin{exercise}{103}{常规·三阶等价量}
已知 $\tan x-x\sim x^3/3$（$x\to0$），求 $\lim_{x\to0}(\tan2x-2x)/x^3$。
\end{exercise}
\begin{exercise}{104}{基础·乘积等价量}
求 $\lim_{x\to0}\frac{(2^{x^2}-1)\arctan2x}{x^3}$。
\end{exercise}
\begin{exercise}{105}{基础·等价关系反求参数}
求常数 $a$，使 $1/(ax^2+1)\sim1/(2x^2+1)$（$x\to+\infty$）。
\end{exercise}
\begin{exercise}{106}{常规·无穷小阶数比较}
当 $x\to0$ 时，求 $(\sqrt{1+x^2}-1)\ln(1+x^2)$ 的等价无穷小，并判断它是否为 $x^2$ 的高阶无穷小。
\end{exercise}
\begin{exercise}{107}{常规·乘积中的等价代换}
求 $\lim_{x\to0}\frac{(1-\cos x)(e^x-1)}{x^2\sin x}$。
\end{exercise}
\begin{exercise}{108}{基础·等价关系反求系数}
求常数 $a$，使 $\sin(ax)\sim3x$（$x\to0$）。
\end{exercise}
\begin{exercise}{109}{常规·对数极限反推}
已知 $f(x)\to0$ 且 $\lim_{x\to0}\ln(1+f(x))/x^2=3$。求 $\lim_{x\to0}f(x)/x^2$。
\end{exercise}
\begin{exercise}{110}{基础·已知极限运算}
已知 $\lim_{x\to0}(\sin6x+xf(x))/x=0$。求 $\lim_{x\to0}f(x)$。
\end{exercise}
\begin{exercise}{111}{常规·振荡项与等价量}
求 $\lim_{x\to0}\frac{4\sin x+3x^2\cos(1/x^2)}{(1+\cos x)\ln(1+2x)}$。
\end{exercise}
\begin{exercise}{112}{基础·根式等价量}
求 $\lim_{x\to0}(\sqrt{1+\sin x}-1)/x$。
\end{exercise}
\begin{exercise}{113}{基础·指数复合等价量}
求 $\lim_{x\to0}(e^{1-\cos x}-1)/x^2$。
\end{exercise}
\begin{exercise}{114}{常规·对数差化简}
求 $\lim_{x\to0}\frac{\ln(x^2+e^{2x})-2x}{x^2}$。提示：先把分子写为 $\ln(1+x^2e^{-2x})$。
\end{exercise}
\begin{exercise}{115}{常规·无穷远换元}
求 $\lim_{x\to+\infty}x\sin\ln(1+4/x)$。
\end{exercise}
\begin{exercise}{116}{常规·积分无穷小}
已知 $\int_0^x(e^t-1)\,\dd t=e^x-1-x$。利用 $e^x=1+x+x^2/2+o(x^2)$，求该积分在 $x\to0$ 时的等价无穷小。
\end{exercise}
\begin{exercise}{117}{常规·相消后的极限}
求 $\lim_{x\to0}(1/\sin x-1/x)$。提示：通分，可使用 $x-\sin x\sim x^3/6$。
\end{exercise}
\begin{exercise}{118}{常规·积分比值极限}
求 $\lim_{x\to0}\frac{x^3}{\int_0^x\ln(1+t^2)\,\dd t}$。可使用积分中值定理及 $\ln(1+t^2)\sim t^2$ 的估计，或使用洛必达法则。
\end{exercise}
\begin{exercise}{119}{基础·等价代换变式}
求 $\lim_{x\to0}\frac{\sin2x\,\ln(1+3x)}{x^2}$。
\end{exercise}
\begin{exercise}{120}{常规·二阶等价量综合}
已知 $\ln(1+x)=x-x^2/2+o(x^2)$。求 $\lim_{x\to0}\frac{\ln(1+x)-x}{1-\cos x}$。
\end{exercise}
\subsection{函数的连续性与间断点}
\begin{exercise}{121}{基础·连续函数值}
已知 $f$ 在 $1$ 处连续，且 $\lim_{x\to1}(f(x)-4)/(x-1)=2$。求 $f(1)$。
\end{exercise}
\begin{exercise}{122}{常规·连续性的定义}
用 $\varepsilon$--$\delta$ 定义证明 $f(x)=2x+1$ 在 $x=1$ 处连续。
\end{exercise}
\begin{exercise}{123}{基础·连续补定义}
设 $f(x)=(e^{2x}-1)/\sin4x$（$x\ne0$），$f(0)=k$，在 $0$ 的充分小邻域内考虑。求使 $f$ 在 $0$ 处连续的 $k$。
\end{exercise}
\begin{exercise}{124}{基础·分段连续性}
设 $f(x)=\begin{cases}\sin2x/\tan x,&x>0,\\2,&x=0,\\\ln(1-2x),&x<0.\end{cases}$ 各支在 $0$ 的充分小邻域内考虑。求左右极限，判断 $f$ 在 $0$ 处是否连续。
\end{exercise}
\begin{exercise}{125}{常规·振荡因子的连续性}
设 $f(x)=\begin{cases}(x-1)^2\cos(1/(x-1)),&x<1,\\(x+1)\ln x,&x\ge1.\end{cases}$ 判断 $f$ 在 $1$ 处是否连续。
\end{exercise}
\begin{exercise}{126}{常规·根式连续延拓}
为 $f(x)=\frac{\sqrt{1+x^2}-\sqrt{1-x^2}}{x\arcsin x}$（$0<|x|\le1$）补定义 $f(0)$，使其在 $0$ 处连续。
\end{exercise}
\begin{exercise}{127}{常规·数列极限生成函数}
对 $x>0$，设 $f(x)=\lim_{n\to\infty}\frac1n\ln(e^{nx}+x^2)$。求 $f(x)$。提示：将 $e^{nx}$ 提出。
\end{exercise}
\begin{exercise}{128}{常规·连续函数值反推}
设 $f$ 在 $0$ 处连续，且 $\lim_{x\to0}\frac{\sin(xf(x))}{x}=2$。求 $f(0)$。提示：可使用 $\sin u=u+o(u)$。
\end{exercise}
\begin{exercise}{129}{基础·有理函数间断点}
求 $f(x)=(x^2-1)/(x^2-3x+2)$ 的间断点并分类，对可去间断点作连续补定义。
\end{exercise}
\begin{exercise}{130}{常规·指数函数间断点}
求 $f(x)=1/(e^x-1)$ 的间断点及该点的左右极限，并分类。
\end{exercise}
\begin{exercise}{131}{常规·连续性反求参数}
设 $f(x)=\begin{cases}x\sin(1/x),&x<0,\\b,&x=0,\\a+\ln(1+3x)/x,&x>0.\end{cases}$ 求使 $f$ 在 $0$ 处连续的 $a,b$。
\end{exercise}
\begin{exercise}{132}{常规·点态极限函数}
设 $f(x)=\lim_{n\to\infty}1/(1+x^{2n})$。求分段表达式，并判断 $x=\pm1$ 处的连续性。
\end{exercise}
\begin{exercise}{133}{常规·夹逼生成函数}
利用 $\max\{1,x\}\le(1+x^n)^{1/n}\le2^{1/n}\max\{1,x\}$（$x\ge0$），求 $f(x)=\lim_{n\to\infty}(1+x^n)^{1/n}$，判断其连续性。
\end{exercise}
\begin{exercise}{134}{常规·分段连续性参数}
设 $f(x)=\begin{cases}\ln(1+2x)/x,&x<0,\\a,&x=0,\\(e^{bx}-1)/x,&x>0.\end{cases}$ 在 $0$ 的充分小邻域内考虑。求使 $f$ 在 $0$ 处连续的 $a,b$。
\end{exercise}
\begin{exercise}{135}{基础·连续函数运算}
求 $f(x)=e^{x/\sin x}$ 在 $x=0$ 处的极限，并作连续补定义。
\end{exercise}
\begin{exercise}{136}{基础·有理无理分段函数}
设 $f(x)=x^2$（$x\in\mathbb Q$），$f(x)=0$（$x\notin\mathbb Q$）。利用 $0\le f(x)\le x^2$ 判断 $f$ 在 $0$ 处是否连续。
\end{exercise}
\begin{exercise}{137}{常规·振荡函数连续补定义}
设 $f(x)=x^2\sin(1/x)$（$x\ne0$），$f(0)=c$。求使 $f$ 在 $0$ 处连续的 $c$，并判断 $\sin(1/x)$ 在 $0$ 处能否连续补定义。
\end{exercise}
\subsection{连续函数的运算与初等函数的连续性}
\begin{exercise}{138}{基础·连续函数商}
求 $\lim_{x\to0}\lg(100+x^2)/(1+\arcsin2x)$。
\end{exercise}
\begin{exercise}{139}{常规·连续函数复合}
设 $f$ 在 $0$ 处连续且 $f(0)=1$。求 $\lim_{x\to0}e^{f(x)}$、$\lim_{x\to0}\ln(1+f(x)^2)$。
\end{exercise}
\begin{exercise}{140}{基础·反函数连续性}
设 $f(x)=x^3$，求其反函数，说明反函数在实轴上连续。
\end{exercise}
\begin{exercise}{141}{常规·连续运算综合}
设 $f$ 在 $[a,b]$ 上连续且 $f(x)\ge1$。说明 $|f|$、$1/f$ 在 $[a,b]$ 上连续，并证明 $1/f$ 有界。
\end{exercise}
\subsection{闭区间上连续函数的性质}
\begin{exercise}{142}{常规·零点存在性}
证明方程 $x^3-3x+1=0$ 在 $(0,1)$ 内至少有一个实根。提示：对 $F(x)=x^3-3x+1$ 使用零点定理。
\end{exercise}
\begin{exercise}{143}{常规·零点存在与唯一性}
证明方程 $x=\frac12\cos x$ 在 $[0,1]$ 上有唯一实根。提示：考察 $F(x)=x-\frac12\cos x$ 的端点符号及单调性。
\end{exercise}
\begin{exercise}{144}{常规·平移差与零点}
设 $f$ 在 $[0,1]$ 上连续且 $f(0)=f(1)$。证明存在 $\xi\in[0,1/2]$ 使 $f(\xi+1/2)=f(\xi)$。提示：考察 $F(x)=f(x+1/2)-f(x)$ 的端点值。
\end{exercise}
\begin{exercise}{145}{常规·连续函数有界性}
设 $f$ 在 $[0,+\infty)$ 上连续且 $\lim_{x\to+\infty}f(x)=1$。证明 $f$ 有界。提示：将定义域分成一个闭区间和一个尾部区间。
\end{exercise}
\begin{exercise}{146}{常规·介值定理}
设 $f$ 在 $[a,b]$ 上连续。证明存在 $\xi\in[a,b]$，使 $f(\xi)=(f(a)+f(b))/2$。
\end{exercise}
\begin{exercise}{147}{基础·最值定理}
求 $f(x)=x(2-x)$ 在 $[0,2]$ 上的最大值与最小值，并说明最值存在的依据。
\end{exercise}
\begin{exercise}{148}{常规·零点定理应用}
设 $f$ 在 $[0,1]$ 上连续，$f(0)=-1$、$f(1)=2$。证明存在 $\xi\in(0,1)$，使 $f(\xi)=\xi$。提示：考虑 $F(x)=f(x)-x$。
\end{exercise}
\begin{exercise}{149}{基础·零点定理变式}
证明方程 $e^x+x-3=0$ 在 $(0,1)$ 内至少有一个实根。
\end{exercise}
\begin{exercise}{150}{常规·闭区间性质综合}
设 $f$ 在 $[a,b]$ 上连续且处处为正。用最小值定理证明存在 $m>0$，使 $f(x)\ge m$，进而说明 $1/f$ 在 $[a,b]$ 上有界。
\end{exercise}

\clearpage
\section{第2章 导数与微分}
\subsection{导数的基本概念}
\begin{exercise}{151}{基础·导数定义}
用导数定义求 $f(x)=1+2x+3x^2$ 在 $x=1$ 处的导数。
\end{exercise}
\begin{exercise}{152}{基础·零点处求导}
不展开多项式，求 $f(x)=x(2x+1)(3x+2)(4x+3)$ 的 $f'(0)$。
\end{exercise}
\begin{exercise}{153}{常规·定义求导与切线}
用导数定义求 $f(x)=x(x-2)\sqrt{\frac{3x^2+1}{2x^2+1}}$ 的 $f'(0)$，并求曲线在原点的切线。
\end{exercise}
\begin{exercise}{154}{基础·补定义函数的导数}
设 $f(x)=\frac{\sin^2x}{2x}$（$x\ne0$），$f(0)=0$。用定义求 $f'(0)$。
\end{exercise}
\begin{exercise}{155}{常规·连续因子与求导}
设 $f(x)=(x^2-1)\varphi(x-1)$，$\varphi$ 在 $0$ 处连续且 $\varphi(0)=3$。用定义求 $f'(1)$。
\end{exercise}
\begin{exercise}{156}{常规·平移差商}
设 $f$ 在 $0$ 处可导，$f'(0)=\frac{1}{2}$，且 $f(2+x)=2f(x)$ 对所有实数 $x$ 成立。用定义求 $f'(2)$。
\end{exercise}
\begin{exercise}{157}{基础·差商极限}
设 $f$ 在 $0$ 处连续，且 $\lim_{x\to0}\frac{f(x)+3}{x}=2$。求 $f(0)$ 和 $f'(0)$。
\end{exercise}
\begin{exercise}{158}{基础·增量缩放}
设 $f'(x_0)=2$。求 $\lim_{h\to0}\frac{f(x_0+5h)-f(x_0)}{h}$ 及 $\lim_{h\to0}\frac{f(x_0)-f(x_0-3h)}{h}$。
\end{exercise}
\begin{exercise}{159}{常规·双移位差商}
设 $f'(x_0)=3$。求 $\lim_{h\to0}\frac{f(x_0+4h)-f(x_0+2h)}{h}$。提示：在分子中加减 $f(x_0)$。
\end{exercise}
\begin{exercise}{160}{常规·二次增量差商}
设 $f$ 在 $3$ 处可导，且 $\lim_{x\to0}\frac{f(3)-f(3-3x^2)}{\ln(1+x^2)}=3$。求 $f'(3)$。
\end{exercise}
\begin{exercise}{161}{基础·已知导数求极限}
设 $f(0)=0$、$f'(0)=1$。求 $\lim_{x\to0}\frac{f(3x)-f(-5x)}{x}$。
\end{exercise}
\begin{exercise}{162}{基础·导数定义换元}
设 $f(0)=0$、$f'(0)=2$。求 $\lim_{h\to0}\frac{f(h^3)}{h^3}$ 和 $\lim_{h\to0}\frac{f(h^2)}{h^2}$。
\end{exercise}
\begin{exercise}{163}{常规·复合增量极限}
设 $f$ 在 $1$ 处可导，$f(1)=0$、$f'(1)=2$。求 $\lim_{x\to0}\frac{f(\cos x)}{x^2}$。
\end{exercise}
\begin{exercise}{164}{基础·左右导数}
分别用左右差商判断 $|x|$、$x^2$ 在 $0$ 处是否可导，可导时求导数。
\end{exercise}
\begin{exercise}{165}{基础·分段函数可导性}
设 $f(x)=\cos x$（$x>0$），$f(x)=x^2+1$（$x\le0$）。求 $f'_-(0)$、$f'_+(0)$，判断可导性。
\end{exercise}
\begin{exercise}{166}{常规·指数分段函数}
设 $f(x)=x(1+e^{\frac{1}{x}})$（$x\ne0$），$f(0)=0$。判断 $f$ 在 $0$ 处的连续性，并求左导数（若存在）。
\end{exercise}
\begin{exercise}{167}{基础·绝对值尖点}
求 $f(x)=x^2|x-1|$ 在 $1$ 处的左右导数，判断是否可导。
\end{exercise}
\begin{exercise}{168}{常规·零点消除尖点}
用定义求 $f(x)=(x^3-1)|x-1|$ 在 $1$ 处的导数。
\end{exercise}
\begin{exercise}{169}{常规·连续因子差商}
设 $\varphi$ 在 $1$ 处连续且 $\varphi(1)=0$，$f(x)=|x-1|\varphi(x)$。用 $|\varphi(x)|\to0$ 证明 $f'(1)=0$。
\end{exercise}
\begin{exercise}{170}{常规·绝对值复合}
求 $f(x)=|x^2-1|$ 的不可导点，并计算这些点的左右导数。
\end{exercise}
\begin{exercise}{171}{基础·切线与法线}
求曲线 $y=2x+\cos x$ 在 $(0,1)$ 处的切线与法线方程。
\end{exercise}
\begin{exercise}{172}{常规·切线垂直}
求实数 $x_0$，使曲线 $y=x^2$ 与 $y=x^3$ 在横坐标同为 $x_0$ 的点处的切线互相垂直。
\end{exercise}
\begin{exercise}{173}{常规·指定方向的切线}
求与直线 $2x-6y+1=0$ 垂直且与曲线 $y=x^3+3x^2-5$ 相切的直线方程。
\end{exercise}
\begin{exercise}{174}{常规·切线与面积}
求双曲线 $xy=1$ 在 $(1,1)$ 处的切线，并求该切线与坐标轴围成的三角形面积。
\end{exercise}
\begin{exercise}{175}{常规·公共切线}
求 $a,b$，使 $y=1+\ln x$ 与 $y=x^2+ax+b$ 在 $x=1$ 处具有同一条切线。
\end{exercise}
\begin{exercise}{176}{常规·偶函数导数}
设 $f$ 为可导偶函数且 $f'(3)=-\frac{1}{2}$。求曲线在 $(-3,f(-3))$ 处的切线斜率。可使用偶函数的导函数为奇函数。
\end{exercise}
\begin{exercise}{177}{基础·可导与连续}
设 $f$ 在 $2$ 处可导且 $\lim_{x\to2}f(x)=5$。求 $f(2)$，并说明依据。
\end{exercise}
\begin{exercise}{178}{常规·振荡函数可导性}
设 $f(x)=x^2\cos(\frac{1}{x})$（$x\ne0$），$f(0)=0$。用定义判断其在 $0$ 处的连续性与可导性。
\end{exercise}
\begin{exercise}{179}{常规·分段拼接}
设 $f(x)=xe^x$（$x\ge0$），$f(x)=x^3+x$（$x<0$）。判断其在 $0$ 处是否可导，并求 $f'(0)$。
\end{exercise}
\begin{exercise}{180}{基础·先连续再可导}
设 $f(x)=x^3+\ln x$（$x\ge1$），$f(x)=x^3\ln x$（$0<x<1$）。判断 $f$ 在 $1$ 处是否连续、可导。
\end{exercise}
\begin{exercise}{181}{常规·可导拼接参数}
求 $a,b$，使 $f(x)=x\sin\pi x$（$x>1$），$f(x)=ax^2+b$（$x\le1$）在 $1$ 处可导。
\end{exercise}
\begin{exercise}{182}{常规·有界因子与可导}
设 $\varphi$ 在 $0$ 附近有界，$f(x)=x^2\varphi(x)$ 且 $f(0)=0$。用定义求 $f'(0)$。
\end{exercise}
\begin{exercise}{183}{基础·导数定义变式}
用定义求 $f(x)=\sin2x$ 的 $f'(0)$。
\end{exercise}
\begin{exercise}{184}{常规·分段可导综合}
设 $f(x)=ax+b$（$x<1$），$f(x)=x^2$（$x\ge1$）。求使 $f$ 在 $1$ 处可导的 $a,b$，并写出切线方程。
\end{exercise}
\subsection{函数的求导法则}
\begin{exercise}{185}{基础·基本求导}
求 $f(x)=3x^2+2e^x-5\ln x+\sin1$ 的导数及 $f'(1)$。
\end{exercise}
\begin{exercise}{186}{基础·乘积法则}
求 $f(x)=x\cos x+3x^2$ 的导数及 $f'(\frac{\pi}{2})$。
\end{exercise}
\begin{exercise}{187}{基础·指数三角乘积}
求 $y=e^x(\cos x-2\sin x)$ 的导数。
\end{exercise}
\begin{exercise}{188}{基础·乘积与复合求导}
求 $y=(1+x^2)\arctan x-x\ln x$（$x>0$）的导数。
\end{exercise}
\begin{exercise}{189}{基础·商式求导}
求 $y=\frac{\sqrt x+xe^x-3x^3\sin x}{x^2}$（$x>0$）的导数。提示：可先拆分为三项。
\end{exercise}
\begin{exercise}{190}{基础·商法则}
求 $y=\frac{\sin x}{x-\cos x}$ 的导数，注明分母限制。
\end{exercise}
\begin{exercise}{191}{基础·对数有理式}
求 $y=\frac{2\ln x+x^3}{3\ln x+x^2}$ 的导数及 $y'(1)$。
\end{exercise}
\begin{exercise}{192}{基础·乘商求导}
求 $y=\frac{e^x(x^2+1)}{x+1}$（$x\ne-1$）的导数。
\end{exercise}
\begin{exercise}{193}{常规·换元与求导}
已知 $f(x-1)=x^2+1$ 对所有实数 $x$ 成立。求 $f(t)$ 和 $f'(t)$。
\end{exercise}
\begin{exercise}{194}{常规·分段导函数}
设 $f(x)=\frac{x}{2x-1}$（$x<0$），$f(x)=e^x\sin x$（$x\ge0$）。分别求 $x<0$、$x>0$ 时的导数，并判断 $0$ 处是否可导。
\end{exercise}
\begin{exercise}{195}{基础·反函数导数}
设 $f$ 的可导反函数为 $\varphi$，已知 $f(1)=2$、$f(2)=3$、$f'(1)=4$、$f'(2)=5$。求 $\varphi'(2)$、$\varphi'(3)$。
\end{exercise}
\begin{exercise}{196}{常规·反函数定点导数}
设 $f(x)=3x^3+2x+1$ 的反函数为 $\varphi$。求 $\varphi'(1)$。提示：先求 $f(x)=1$ 的解。
\end{exercise}
\begin{exercise}{197}{基础·链式法则}
求 $y=(1-x+2x^3)^4$ 的导数及 $y'(0)$。
\end{exercise}
\begin{exercise}{198}{基础·多层复合求导}
求 $y=\ln(\cos(e^{-x}))$ 在 $x=0$ 附近的导数及 $y'(0)$。
\end{exercise}
\begin{exercise}{199}{常规·反三角复合求导}
求 $f(x)=\arcsin\sqrt{\frac{1-x}{1+x}}$ 在 $0<x<1$ 内的导数及 $f'(\frac{1}{2})$。
\end{exercise}
\begin{exercise}{200}{常规·根式商求导}
求 $y=\frac{\sqrt{x+1}-\sqrt{x-1}}{\sqrt{x+1}+\sqrt{x-1}}$（$x>1$）的导数。提示：可先化简为 $x-\sqrt{x^2-1}$。
\end{exercise}
\begin{exercise}{201}{基础·指数与幂指求导}
求 $y=3^x+x^3+x^{2x}$（$x>0$）的导数及 $y'(1)$。
\end{exercise}
\begin{exercise}{202}{基础·三角复合求导}
求 $y=\cos3x\,\sin^3x$ 的导数及 $y'(\frac{\pi}{6})$。
\end{exercise}
\begin{exercise}{203}{基础·反三角与根式求导}
求 $y=x\sqrt{1-x^2}-\arccos x$ 在 $-1<x<1$ 内的导数。
\end{exercise}
\begin{exercise}{204}{常规·根式对数求导}
求 $y=x^2+x\sqrt{x^2+1}+\ln(x+\sqrt{x^2+1})$ 的导数并化简。
\end{exercise}
\begin{exercise}{205}{基础·对数三角复合}
求 $y=x[\cos(\ln x)+\sin(\ln x)]$（$x>0$）的导数。
\end{exercise}
\begin{exercise}{206}{基础·反余切复合}
求 $y=e^{-x}\operatorname{arccot}(e^x)$ 的导数，其中 $\operatorname{arccot}\in(0,\pi)$。
\end{exercise}
\begin{exercise}{207}{常规·先求极限再求导}
对固定实数 $x$，设 $f(x)=\lim_{n\to\infty}x(1-\frac{3x}{n})^n$，只考虑底数为正的充分大的 $n$。求 $f(x)$ 和 $f'(x)$。
\end{exercise}
\begin{exercise}{208}{基础·抽象函数链式法则}
设 $f$ 可导，求 $y=f(3^x+x^2)$ 的导数；若 $f'(1)=2$，求 $y'(0)$。
\end{exercise}
\begin{exercise}{209}{常规·复合函数取点}
已知 $f$ 可导且 $f'(1)=-2$，$g(x)=f(2x^3-3x)$。求 $g'(-1)$。
\end{exercise}
\begin{exercise}{210}{基础·复合求导变式}
求 $y=\ln(1+x^2)$ 的导数及在 $x=1$ 处的切线。
\end{exercise}
\begin{exercise}{211}{常规·反函数复合求导}
设 $f(x)=x+e^x$ 的反函数为 $\varphi$，$G(x)=\varphi(e^x)$。求 $\varphi'(1)$ 和 $G'(0)$。
\end{exercise}
\subsection{高阶导数}
\begin{exercise}{212}{基础·二阶导数}
分别求 $y=\sin2x$ 和 $y=x\arctan x$ 的二阶导数。
\end{exercise}
\begin{exercise}{213}{基础·指数三角二阶导数}
设 $f(x)=e^{3x}\sin2x$。求 $f''(x)$ 和 $f''(\frac{\pi}{2})$。
\end{exercise}
\begin{exercise}{214}{基础·对数函数二阶导数}
设 $y=\frac12\ln(1+e^{2x})-x$。求 $y'(0)$ 和 $y''(0)$。
\end{exercise}
\begin{exercise}{215}{基础·反正弦二阶导数}
设 $y=\arcsin(\frac{x}{3})$。求 $y''(x)$（$|x|<3$）及 $y''(\sqrt5)$。
\end{exercise}
\begin{exercise}{216}{常规·高阶导数通式}
设 $y=x\ln x$（$x>0$）。求 $y''$、$y'''$，并归纳写出 $y^{(n)}$（$n\ge2$）的通式。
\end{exercise}
\begin{exercise}{217}{基础·三角高阶导数}
设 $y=(\sin(\frac{x}{2})+\cos(\frac{x}{2}))^2$。先化简，再求 $y^{(10)}$。
\end{exercise}
\begin{exercise}{218}{常规·莱布尼茨公式}
设 $f(x)=xe^{-x}$。用莱布尼茨公式求 $f^{(n)}(x)$（$n\ge1$）。
\end{exercise}
\begin{exercise}{219}{基础·多项式高阶导数}
设 $y=(x^2+1)^5(x^5+1)^2$。不完全展开，求 $y^{(20)}$ 和 $y^{(21)}$。
\end{exercise}
\begin{exercise}{220}{常规·降幂求导}
设 $f(x)=\sin^4x+\cos^4x$。利用降幂公式求 $f^{(4)}(x)$。
\end{exercise}
\begin{exercise}{221}{常规·对数高阶导数}
设 $f(x)=\ln(2-3x)$（$x<\frac{2}{3}$）。求 $f'$、$f''$、$f'''$，并写出 $f^{(n)}$ 的通式。
\end{exercise}
\begin{exercise}{222}{常规·导数关系再求导}
设 $f'=f^2$ 且 $f(1)=2$。求 $f''(1)$。已知 $f$ 在 $1$ 附近可导。
\end{exercise}
\begin{exercise}{223}{基础·高阶导数变式}
求 $y=e^{2x}$ 的四阶导数及 $y^{(4)}(0)$。
\end{exercise}
\begin{exercise}{224}{常规·乘积高阶导数综合}
用莱布尼茨公式求 $f(x)=x^2e^x$ 的 $f^{(4)}(x)$，并求 $f^{(4)}(0)$。
\end{exercise}
\subsection{隐函数及由参数方程所确定的函数的导数}
\begin{exercise}{225}{基础·隐函数求导}
由方程 $x^2+y^2=6xy$ 求 $\frac{\dd y}{\dd x}$，注明公式分母非零的条件。
\end{exercise}
\begin{exercise}{226}{常规·隐函数定点导数}
方程 $xy-\sin(\pi y^2)=0$ 确定过 $(0,1)$ 的可导分支。求 $y'(0)$。
\end{exercise}
\begin{exercise}{227}{常规·隐曲线切线}
求曲线 $\tan(x+y+\frac{\pi}{4})=y$ 在 $(-1,1)$ 处的切线与法线。
\end{exercise}
\begin{exercise}{228}{常规·隐函数分支}
求曲线 $y^2+2e^{xy}=3y$ 在 $(0,1)$ 处的切线方程。
\end{exercise}
\begin{exercise}{229}{常规·抽象隐函数求导}
设 $xf(y)=e^y\ln y$ 确定过 $(0,1)$ 的可导分支，$f$ 可导且 $f(1)=1$。求 $y'(0)$。
\end{exercise}
\begin{exercise}{230}{常规·公共切线参数}
曲线 $y=x^2+ax+b$ 与 $2y=xy^2-3$ 在 $(1,-1)$ 处有公共切线。求 $a,b$。
\end{exercise}
\begin{exercise}{231}{常规·隐函数二阶导数}
方程 $x^2+y^2=2$ 确定过 $(1,1)$ 的分支。求 $y'(1)$ 和 $y''(1)$。
\end{exercise}
\begin{exercise}{232}{常规·指数隐函数二阶导数}
设 $e^y+xy=e+1$ 确定过 $(1,1)$ 的局部分支。求 $y'(1)$、$y''(1)$。
\end{exercise}
\begin{exercise}{233}{常规·含导数的函数关系}
设 $y(x)$ 二阶可导，经过 $(1,-1)$，满足 $x^2+y^2+(y')^3=1$。求 $y'(1)$、$y''(1)$。
\end{exercise}
\begin{exercise}{234}{基础·对数求导}
用对数求导法求 $y=x^{\ln x}$（$x>0$）的导数。
\end{exercise}
\begin{exercise}{235}{基础·变底变指数求导}
设 $y=(1+\sin x)^{x^2}$，其中 $1+\sin x>0$。求 $y'$ 和 $y'(0)$。
\end{exercise}
\begin{exercise}{236}{常规·幂指函数乘积}
求 $y=x^{\sin x}\sqrt{e^x-1}$（$x>0$）的导数。
\end{exercise}
\begin{exercise}{237}{常规·复合幂指函数}
求 $y=(1+2x)^{\frac{1}{x}}e^{-x}$（$x>-\frac{1}{2}$、$x\ne0$）的导数。
\end{exercise}
\begin{exercise}{238}{常规·多因子对数求导}
用对数求导法求 $y=\frac{(x+1)^2\sqrt{2x+2}}{(x-1)^3}$（$x>1$）的导数。
\end{exercise}
\begin{exercise}{239}{基础·幂指函数}
求 $y=x^x$（$x>0$）的导数及驻点。
\end{exercise}
\begin{exercise}{240}{常规·幂指二阶导数}
求 $y=x^{2x}$（$x>0$）的二阶导数。
\end{exercise}
\begin{exercise}{241}{常规·对数化隐函数}
在 $x,y>0$ 且 $1+\ln y\ne0$ 的可导分支上，方程 $x^{\frac{1}{y}}=y^{\frac{1}{x}}$ 确定 $y=y(x)$。取对数化简后求 $y'$。
\end{exercise}
\begin{exercise}{242}{基础·参数方程求导}
求参数曲线 $x=\frac{2t}{1+t^2}$、$y=\frac{2t^2}{1+t^2}$ 的 $\frac{\dd y}{\dd x}$，注明 $\frac{\dd x}{\dd t}\ne0$ 的限制。
\end{exercise}
\begin{exercise}{243}{基础·参数定点导数}
设 $x=t^2e^{2t}$、$y=(t^2+1)e^t$。求 $t=2$ 时的 $\frac{\dd y}{\dd x}$。
\end{exercise}
\begin{exercise}{244}{常规·参数曲线切线}
求参数曲线 $x=\sqrt{t^2+3}$、$y=\ln(t+\sqrt{t^2+3})$ 在 $t=1$ 对应点处的切线。
\end{exercise}
\begin{exercise}{245}{常规·参数曲线法线}
求参数曲线 $x=e^t\sin2t$、$y=e^t\cos t$ 在 $t=0$ 对应点处的切线和法线。
\end{exercise}
\begin{exercise}{246}{常规·参数二阶导数}
设 $x=\cos t$、$y=t\cos t-\sin t$，$0<t<\pi$。求 $\frac{\dd y}{\dd x}$ 与 $\frac{\dd^2y}{\dd x^2}$。
\end{exercise}
\begin{exercise}{247}{常规·参数二阶导数取值}
设 $x=\arctan t$、$y=3t+t^3$。求 $t=1$ 时的 $\frac{\dd^2y}{\dd x^2}$。
\end{exercise}
\begin{exercise}{248}{常规·参数导数综合}
设 $x=e^t+t$、$y=\sin t$。求 $t=0$ 时的 $\frac{\dd y}{\dd x}$ 与 $\frac{\dd^2y}{\dd x^2}$。
\end{exercise}
\begin{exercise}{249}{基础·隐函数求导变式}
求曲线 $x^2+xy+y^2=3$ 在 $(1,1)$ 处的导数与切线。
\end{exercise}
\begin{exercise}{250}{常规·参数曲线小综合}
设 $x=t^2+1$、$y=t^3$。求 $t=1$ 对应点的切线及该点处的二阶导数。
\end{exercise}
\subsection{函数的微分}
\begin{exercise}{251}{基础·微分与增量}
设 $y=x^3+x$。在 $x=3$、$\Delta x=0.1$ 时计算 $\dd y$ 和 $\Delta y$。
\end{exercise}
\begin{exercise}{252}{基础·由增量求微分}
已知 $f$ 的增量为 $\Delta y=\frac{\Delta x}{\sqrt{1+x^3}}+o(\Delta x)$（$x>-1$）。求 $f'(2)$ 及 $x=2$ 处的微分。
\end{exercise}
\begin{exercise}{253}{常规·微分的线性主部}
设 $f$ 在 $x_0$ 处可导且 $f'(x_0)=1$。写出 $\dd y$，并求 $\lim_{\Delta x\to0}\frac{\Delta y}{\Delta x}$。
\end{exercise}
\begin{exercise}{254}{基础·复合函数微分}
求 $y=e^{\sin^2x}$ 和 $y=\arcsin(\frac{x}{2})$（$|x|<2$）的微分。
\end{exercise}
\begin{exercise}{255}{基础·定点微分}
分别求 $y=\sqrt2+\sin x$ 在 $x=0$ 处、$y=\ln(1+2^{-2x})$ 在 $x=1$ 处的微分。
\end{exercise}
\begin{exercise}{256}{常规·隐函数微分}
设 $2y+\sin(x+y)=0$ 确定可导函数 $y=y(x)$。求 $\dd y$ 及原点处的微分。
\end{exercise}
\begin{exercise}{257}{常规·隐函数定点微分}
方程 $e^y+xy-e=0$ 在 $(0,1)$ 附近确定分支。求 $x=0$ 处的微分，并估计 $x$ 增加 $0.01$ 时 $y$ 的增量。
\end{exercise}
\begin{exercise}{258}{常规·补定义点微分}
设 $f(x)=\frac{\cos x-1}{x}$（$x\ne0$），$f(0)=0$。用导数定义求 $x=0$ 处的微分。
\end{exercise}
\begin{exercise}{259}{基础·幂指函数微分}
求 $y=(1+e^x)^x$ 的微分及 $x=0$ 处的微分。
\end{exercise}
\begin{exercise}{260}{基础·由导数求函数}
求导数分别为 $x^2$、$\cos3x$、$e^x$ 的全部函数，各式均在实轴上考虑。
\end{exercise}
\begin{exercise}{261}{常规·夹逼与可微性}
设 $f$ 在 $0$ 附近有定义且 $|f(x)|\le x^2$。用导数定义证明 $f'(0)=0$，并写出该点的微分。
\end{exercise}
\begin{exercise}{262}{常规·定点导数计算}
分别求 $p(x)=x(x-1)(x-2)(x-3)$ 的 $p'(3)$ 和 $y=\frac{(x^2+1)\sqrt{3x+1}}{(x+1)^2e^{2x}}$ 的 $y'(1)$。
\end{exercise}
\begin{exercise}{263}{常规·差商反求导数}
设 $f$ 在 $x_0$ 处可导，$f(x_0)=0$，且 $\lim_{h\to0}\frac{f(x_0+3h)-3f(x_0-h)}{h}=1$。求 $f'(x_0)$。
\end{exercise}
\begin{exercise}{264}{常规·可导信息与幂指极限}
设 $f$ 在 $0$ 处可导，$f(0)=1$、$f'(0)=4$。求 $\lim_{x\to0}[f(\sin x)]^{\frac{1}{2x}}$。提示：取对数并使用导数定义。
\end{exercise}
\begin{exercise}{265}{常规·切线截距}
求曲线 $y=\frac{1}{1+x^2}$ 在 $(1,\frac{1}{2})$ 处的切线及它与 $x$ 轴的交点。
\end{exercise}
\begin{exercise}{266}{常规·复合差商与法线}
设 $f$ 连续可导，且 $\lim_{x\to1}\frac{f(3x-1)+3}{x-1}=3$。求 $f(2)$、$f'(2)$ 及该点的法线方程。
\end{exercise}
\begin{exercise}{267}{常规·分段函数微分}
设 $f(x)=x^2$（$|x|\le1$），$f(x)=1$（$|x|>1$）。判断 $x=1$ 处是否可微，并求 $x=0$ 处的微分。
\end{exercise}
\begin{exercise}{268}{常规·连续补值}
设 $f$ 在 $1$ 处可导，$f(1)=0$、$f'(1)=1$。令 $g(x)=\frac{f(e^{2x})+\ln(1-4x)}{x}$（$x\ne0$），$g(0)=b$，在 $0$ 附近考虑。求使 $g$ 连续的 $b$。
\end{exercise}
\begin{exercise}{269}{常规·分段可导参数}
设 $f(x)=x^2$（$x\ge1$），$f(x)=3ax+b$（$x<1$）。求使 $f$ 在 $1$ 处可导的 $a,b$。
\end{exercise}
\begin{exercise}{270}{基础·互补复合求导}
设 $f$ 可导，$y=f(\sin^2x)+f(\cos^2x)$。求 $y'$ 与 $y'(0)$。
\end{exercise}
\begin{exercise}{271}{常规·两层抽象复合}
设 $f$ 可导，$y=f(f(1+3x^2))$。求 $y'$ 与 $y'(0)$。
\end{exercise}
\begin{exercise}{272}{常规·公共切线与复合导数}
曲线 $y=f(x)$ 与 $y=\frac{3+\ln x}{x}$ 在 $x=1$ 处切线相同。设 $F(x)=[f(e^x)]^2$，求 $F'(0)$。
\end{exercise}
\begin{exercise}{273}{常规·二阶导数与极限}
设 $f$ 在 $0$ 附近二阶连续可导，$f(0)=f'(0)=0$、$f''(0)=2$。求 $\lim_{x\to0}\frac{e^{f(x)}-1}{1-\cos x}$。提示：使用二阶泰勒展开。
\end{exercise}
\begin{exercise}{274}{常规·复合二阶导数}
设 $f$ 二阶可导且 $f'(0)=3$，$y=f(x^2)$。求 $y''(0)$。
\end{exercise}
\begin{exercise}{275}{常规·导数关系计算}
设 $f'=e^f$ 且 $f(0)=0$。求 $f'(0)$ 和 $f''(0)$。
\end{exercise}
\begin{exercise}{276}{常规·莱布尼茨定点求导}
设 $f(x)=(x+1)^3e^{-x^2}$。用莱布尼茨公式求 $f'''(-1)$。
\end{exercise}
\begin{exercise}{277}{常规·反函数二阶导数}
设 $x=\varphi(y)$ 是 $y=x+\ln x$（$x>0$）的反函数。求 $y=1$ 时的 $\frac{\dd x}{\dd y}$ 与 $\frac{\dd^2x}{\dd y^2}$。
\end{exercise}
\begin{exercise}{278}{常规·振荡函数定点求导}
设 $f(x)=x^3\sin(\frac{1}{x})+\sin x+1$（$x\ne0$），$f(0)=1$。用定义求 $f'(0)$，并求 $x\ne0$ 时的导数。
\end{exercise}
\begin{exercise}{279}{常规·隐函数与局部极限}
方程 $\ln(x^2+y)=x^2y+\sin x$ 在 $(0,1)$ 附近确定分支 $y=f(x)$。求 $f'(0)$ 及 $\lim_{h\to0}\frac{f(h)-1}{h}$。
\end{exercise}
\begin{exercise}{280}{常规·隐函数复合求导}
设 $f$ 可导且 $f'(1)=1$，$xy-e^y+e^x=0$ 确定过原点的可导分支。令 $z=f(e^{y(x)})$，求 $z'(0)$。
\end{exercise}
\begin{exercise}{281}{基础·参数定点微分}
设 $x=2t^2+3t+1$、$y=e^t\sin t+2$。求 $t=0$ 时的 $\frac{\dd y}{\dd x}$ 及该点的微分。
\end{exercise}
\begin{exercise}{282}{常规·复合参数求导}
设 $x=t-\frac12\ln(1+e^{2t})$、$y=\arctan(e^t)$。求 $t=0$ 时的 $\frac{\dd y}{\dd x}$。
\end{exercise}
\begin{exercise}{283}{常规·极坐标曲线切线}
求螺线 $\rho=\theta$ 在 $\theta=\frac{\pi}{2}$ 处的切线方程。提示：使用 $x=\theta\cos\theta$、$y=\theta\sin\theta$。
\end{exercise}
\begin{exercise}{284}{基础·由微分求导数}
设 $y=f(e^{2x})$，在 $x=0$ 处有 $\dd y=\dd x$。求 $f'(1)$。
\end{exercise}
\begin{exercise}{285}{常规·复合定点微分}
设 $f'(x)=\arctan(\frac{1}{x})$（$x\ne0$），$y=f(\frac{2x-1}{x+1})$。求 $x=2$ 处的微分。
\end{exercise}
\begin{exercise}{286}{常规·参数函数微分}
设 $f$ 可导且 $f'(0)\ne0$，$x=f(\ln t)$、$y=f(t^2-1)$（$t>0$）。求 $t=1$ 时的 $\dd y$，分别用 $\dd x$ 与 $\dd t$ 表示。
\end{exercise}
\begin{exercise}{287}{基础·微分近似}
用微分估算 $\sqrt{4.04}$。
\end{exercise}
\begin{exercise}{288}{常规·隐函数微分综合}
设 $ye^y=x$ 在 $(e,1)$ 附近确定 $y(x)$。求该点的微分，并用它估算 $x=e+0.02$ 时的 $y$。
\end{exercise}

\clearpage
\section{第3章 微分中值定理与导数的应用}
\subsection{中值定理}
\begin{exercise}{289}{基础·罗尔定理条件}
判断 $f(x)=|x|$ 与 $g(x)=x^2-1$ 在 $[-1,1]$ 上是否满足罗尔定理条件；对满足条件者求罗尔点。
\end{exercise}
\begin{exercise}{290}{基础·罗尔点计算}
验证 $f(x)=\ln\sin x$ 在 $[\frac{\pi}{6},\frac{5\pi}{6}]$ 上满足罗尔定理，并求 $f'(\xi)=0$ 的点。
\end{exercise}
\begin{exercise}{291}{常规·辅助函数与罗尔定理}
设 $f\in C[0,1]$，在 $(0,1)$ 可导，$f(0)=0$、$f(1)=\frac{\pi}{4}$。对 $F(x)=f(x)-\arctan x$ 使用罗尔定理，证明存在 $\xi\in(0,1)$ 使 $(1+\xi^2)f'(\xi)=1$。
\end{exercise}
\begin{exercise}{292}{常规·多项式与罗尔定理}
设 $a+\frac{b}{2}+\frac{c}{3}=0$。对 $F(x)=ax+\frac{bx^2}{2}+\frac{cx^3}{3}$ 使用罗尔定理，证明 $a+bx+cx^2$ 在 $(0,1)$ 内有零点。
\end{exercise}
\begin{exercise}{293}{常规·两次罗尔定理}
设 $f$ 在包含 $[0,1]$ 的开区间二阶可导，$f(0)=f(1)=0$、$f'(0)=0$。证明存在 $\xi\in(0,1)$ 使 $f''(\xi)=0$。提示：先对 $f$ 使用罗尔定理，再对 $f'$ 使用。
\end{exercise}
\begin{exercise}{294}{常规·乘积辅助函数}
设 $f\in C^2[0,1]$ 且 $f(1)=0$，令 $F(x)=x^2f(x)$。求 $F(0)$、$F'(0)$，并用两次罗尔定理证明 $F''$ 在 $(0,1)$ 内有零点。
\end{exercise}
\begin{exercise}{295}{常规·指数辅助函数}
设 $f\in C[0,1]$，在 $(0,1)$ 可导，$f(0)=1$、$f(1)=e^{-1}$。对 $F(x)=f(x)-e^{-x}$ 使用罗尔定理，证明存在 $\xi\in(0,1)$ 使 $f'(\xi)+e^{-\xi}=0$。
\end{exercise}
\begin{exercise}{296}{常规·乘积函数零点}
设 $a>0$，$f\in C[0,a]$，在 $(0,a)$ 可导且 $f(a)=0$。对 $F(x)=xf(x)$ 使用罗尔定理，证明存在 $\xi\in(0,a)$ 使 $f(\xi)+\xi f'(\xi)=0$。
\end{exercise}
\begin{exercise}{297}{常规·加权导数零点}
设 $f\in C[0,1]$，在 $(0,1)$ 可导且 $f(1)=0$。对 $F(x)=x^2f(x)$ 使用罗尔定理，证明存在 $\xi\in(0,1)$ 使 $2f(\xi)+\xi f'(\xi)=0$。
\end{exercise}
\begin{exercise}{298}{基础·对称导数等式}
设 $f$ 在 $(0,1)$ 可导且处处非零。直接取 $\xi=\frac{1}{2}$，验证 $\frac{f'(\xi)}{f(\xi)}=\frac{f'(1-\xi)}{f(1-\xi)}$。
\end{exercise}
\begin{exercise}{299}{常规·平方辅助函数}
设 $f\in C[0,1]$，在 $(0,1)$ 可导，$f(0)=1$、$f(1)=\sqrt e$。对 $F(x)=f(x)^2-e^x$ 使用罗尔定理，证明存在 $\xi\in(0,1)$ 使 $2f(\xi)f'(\xi)=e^\xi$。
\end{exercise}
\begin{exercise}{300}{基础·驻点存在性}
设 $f\in C[0,3]$，在 $(0,3)$ 可导，$f(0)=f(3)$。说明存在 $\xi\in(0,3)$ 使 $f'(\xi)=0$，并写出所用定理。
\end{exercise}
\begin{exercise}{301}{常规·介值与中值定理}
设 $f\in C[0,1]$，在 $(0,1)$ 可导，$f(0)=f(1)=0$、$f(\frac{1}{2})=1$。先对 $f(x)-x$ 使用零点定理，证明存在 $\eta\in(\frac{1}{2},1)$ 使 $f(\eta)=\eta$；再在 $[0,\eta]$ 上证明存在 $\xi$ 使 $f'(\xi)=1$。
\end{exercise}
\begin{exercise}{302}{常规·三零点与二阶导数}
设 $f\in C^2[0,1]$，$f(0)=f(\frac{1}{2})=f(1)=0$。使用两次罗尔定理证明存在 $\xi\in(0,1)$ 使 $f''(\xi)=0$。
\end{exercise}
\begin{exercise}{303}{基础·拉格朗日中值定理}
分别判断 $\ln(1+x^2)$、$|x|$ 在 $[-1,1]$ 上是否满足拉格朗日中值定理条件；对前者求中值点。
\end{exercise}
\begin{exercise}{304}{基础·中值点计算}
对 $f(x)=x^2$ 在 $[1,3]$ 上使用拉格朗日中值定理，求满足 $f(3)-f(1)=2f'(\xi)$ 的 $\xi$。
\end{exercise}
\begin{exercise}{305}{常规·三角函数不等式}
用拉格朗日中值定理证明 $\sin x<x$ 对 $0<x<\frac{\pi}{2}$ 成立。提示：对 $\sin t$ 在 $[0,x]$ 上使用定理。
\end{exercise}
\begin{exercise}{306}{常规·对数增量估计}
设 $a>b>0$。用 $f(x)=\ln x$ 的拉格朗日中值定理证明 $\frac{a-b}{a}<\ln(\frac{a}{b})<\frac{a-b}{b}$。
\end{exercise}
\begin{exercise}{307}{常规·对数不等式}
用拉格朗日中值定理证明 $\frac{x}{1+x}<\ln(1+x)<x$（$x>0$）。提示：对 $\ln(1+t)$ 在 $[0,x]$ 上使用定理。
\end{exercise}
\begin{exercise}{308}{常规·指数不等式}
用拉格朗日中值定理证明 $e^x>1+x$（$x>0$）。
\end{exercise}
\begin{exercise}{309}{常规·导数为零的函数}
在 $(-1,1)$ 上求 $F(x)=\arcsin x+\arccos x$ 的导数，利用中值定理证明 $F(x)=\frac{\pi}{2}$。
\end{exercise}
\begin{exercise}{310}{常规·反三角恒等式}
在 $x>0$ 上用求导证明 $\arctan x+\arctan(\frac{1}{x})=\frac{\pi}{2}$。提示：先证明左侧为常数，再取 $x=1$。
\end{exercise}
\begin{exercise}{311}{常规·乘积函数中值等式}
设 $f\in C[a,b]$，在 $(a,b)$ 可导，$a<b$。对 $F(x)=xf(x)$ 使用拉格朗日中值定理，证明存在 $\xi\in(a,b)$ 使 $bf(b)-af(a)=(b-a)[f(\xi)+\xi f'(\xi)]$。
\end{exercise}
\begin{exercise}{312}{常规·导数上界与增量}
设 $f\in C[a,b]$，在 $(a,b)$ 可导且 $|f'(x)|\le M$。用拉格朗日中值定理证明 $|f(b)-f(a)|\le M(b-a)$。
\end{exercise}
\begin{exercise}{313}{基础·柯西中值点}
对 $f(x)=\sin x$、$g(x)=1+\cos x$ 在 $[0,\frac{\pi}{2}]$ 上使用柯西中值定理，求满足 $[f(\frac{\pi}{2})-f(0)]g'(\xi)=[g(\frac{\pi}{2})-g(0)]f'(\xi)$ 的 $\xi$。
\end{exercise}
\begin{exercise}{314}{常规·柯西中值定理}
设 $0<a<b$，$f\in C[a,b]$，在 $(a,b)$ 可导。对 $F(x)=\frac{f(x)}{x}$、$G(x)=\frac{1}{x}$ 使用柯西中值定理，证明存在 $\xi\in(a,b)$ 使 $\frac{af(b)-bf(a)}{a-b}=f(\xi)-\xi f'(\xi)$。
\end{exercise}
\begin{exercise}{315}{常规·幂函数参照}
设 $0<a<b$，$f\in C[a,b]$，在 $(a,b)$ 可导。对 $f(x)$、$g(x)=x^3$ 使用柯西中值定理，证明存在 $\xi\in(a,b)$ 使 $\frac{f(b)-f(a)}{b-a}=\frac{(a^2+ab+b^2)f'(\xi)}{3\xi^2}$。
\end{exercise}
\begin{exercise}{316}{常规·复合函数中值定理}
对 $F(t)=\sin t$ 在 $[0,1]$ 上使用中值定理，证明存在 $\xi\in(1,e)$ 使 $\sin1=\cos(\ln\xi)$。
\end{exercise}
\begin{exercise}{317}{基础·中值定理变式}
对 $f(x)=e^x$ 在 $[0,1]$ 上使用拉格朗日中值定理，求中值点。
\end{exercise}
\begin{exercise}{318}{常规·罗尔定理综合}
设 $f\in C[0,1]$，在 $(0,1)$ 可导，$f(0)=0$、$f(1)=1$。对 $F(x)=e^x[f(x)-x]$ 使用罗尔定理，证明存在 $\xi\in(0,1)$ 使 $f'(\xi)+f(\xi)-\xi=1$。
\end{exercise}
\subsection{洛必达法则}
\begin{exercise}{319}{基础·零比零型极限}
用洛必达法则求 $\lim_{x\to1}\frac{x^3-1}{x^2-1}$。
\end{exercise}
\begin{exercise}{320}{基础·对数极限}
用洛必达法则求 $\lim_{x\to1}\frac{\ln x}{x-1}$。
\end{exercise}
\begin{exercise}{321}{基础·混合函数极限}
用洛必达法则求 $\lim_{x\to1}\frac{x^3-1+\ln x}{e^x-e}$。
\end{exercise}
\begin{exercise}{322}{基础·无穷比无穷型}
求 $\lim_{x\to+\infty}\frac{x^2}{e^x}$ 与 $\lim_{x\to+\infty}\frac{x^4+e^x}{xe^x}$。
\end{exercise}
\begin{exercise}{323}{常规·连续使用洛必达法则}
求 $\lim_{x\to0}\frac{x-\sin x}{x^2\sin x}$。提示：可先将分母与 $x^3$ 比较，再对 $\frac{x-\sin x}{x^3}$ 使用洛必达法则。
\end{exercise}
\begin{exercise}{324}{基础·对数与幂比较}
求 $\lim_{x\to+\infty}\frac{\ln x}{\sqrt x}$。
\end{exercise}
\begin{exercise}{325}{常规·二阶不定式}
求 $\lim_{x\to0}\frac{e^x-\sin x-1}{1-\sqrt{1-x^2}}$。
\end{exercise}
\begin{exercise}{326}{常规·反正切差极限}
求 $\lim_{x\to0}\frac{\arctan x-x}{x^3}$。
\end{exercise}
\begin{exercise}{327}{基础·乘积极限}
求 $\lim_{x\to0}\frac{\sin x(e^{2x}-1)}{x\ln(1+x)}$。
\end{exercise}
\begin{exercise}{328}{常规·极限反求参数}
求常数 $a$，使 $\lim_{x\to0}\frac{a\tan x+1-\cos x}{\ln(1+2x)}=2$。
\end{exercise}
\begin{exercise}{329}{常规·二阶对数极限}
用洛必达法则求 $\lim_{x\to0}\frac{x-\ln(1+x)}{x^2}$。
\end{exercise}
\begin{exercise}{330}{常规·嵌套三角极限}
求 $\lim_{x\to0}\frac{\sin x-\sin(\sin x)}{x^3}$。提示：令 $u=\sin x$，使用 $u-\sin u\sim\frac{u^3}{6}$。
\end{exercise}
\begin{exercise}{331}{常规·无穷量相减}
求 $\lim_{x\to0}\frac1x(\frac{1}{\sin x}-\frac{1}{\tan x})$。提示：先合并为一个分式。
\end{exercise}
\begin{exercise}{332}{基础·指数衰减}
用洛必达法则求 $\lim_{x\to+\infty}x^3e^{-2x}$。
\end{exercise}
\begin{exercise}{333}{常规·零乘无穷型}
求 $\lim_{x\to1}(1-x)\tan(\frac{\pi x}{2})$。提示：令 $t=1-x$。
\end{exercise}
\begin{exercise}{334}{常规·三角无穷量之差}
求 $\lim_{x\to\frac{\pi}{2}}(\sec x-\tan x)$。提示：乘共轭式或写成 $\frac{\cos x}{1+\sin x}$。
\end{exercise}
\begin{exercise}{335}{常规·倒数差极限}
求 $\lim_{x\to0}[\frac{1}{x}-\frac{1}{e^x-1}]$。提示：先通分。
\end{exercise}
\begin{exercise}{336}{常规·二阶相消}
求 $\lim_{x\to0}[\frac{1}{x^2}-\frac{1}{x\tan x}]$。可使用 $\tan x=x+\frac{x^3}{3}+o(x^3)$。
\end{exercise}
\begin{exercise}{337}{基础·零的零次幂型}
求 $\lim_{x\to0^+}x^x$。提示：先求 $x\ln x$ 的极限。
\end{exercise}
\begin{exercise}{338}{基础·幂指极限}
求 $\lim_{x\to1}x^{\frac{1}{1-x}}$。
\end{exercise}
\begin{exercise}{339}{基础·一的无穷次幂型}
求 $\lim_{x\to0}(1+3x)^{\frac{2}{\sin x}}$，只在底数为正的充分小邻域内考虑。
\end{exercise}
\begin{exercise}{340}{常规·二阶幂指极限}
求 $\lim_{x\to0}(\cos x)^{\frac{1}{\ln(1+x^2)}}$。提示：取对数后使用洛必达法则。
\end{exercise}
\begin{exercise}{341}{基础·无穷的零次幂型}
求 $\lim_{x\to+\infty}x^{\frac{1}{x}}$。
\end{exercise}
\begin{exercise}{342}{常规·指数平均的幂极限}
求 $\lim_{x\to0}[\frac{e^x+e^{3x}}{2}]^{\frac{1}{x}}$。
\end{exercise}
\begin{exercise}{343}{常规·指数差极限}
求 $\lim_{x\to0}\frac{e^{\sin x}-1}{x}$。
\end{exercise}
\begin{exercise}{344}{常规·连续补值参数}
设 $f(x)=\frac{\sin2x+e^{2ax}-1}{x}$（$x\ne0$），$f(0)=a$。求使 $f$ 在 $0$ 处连续的 $a$。
\end{exercise}
\begin{exercise}{345}{基础·幂函数连续补值}
设 $f(x)=(\cos x)^{x^2}$（$0<|x|<\frac{\pi}{2}$），$f(0)=a$。求使 $f$ 在 $0$ 处连续的 $a$。
\end{exercise}
\begin{exercise}{346}{常规·端点相消极限}
求 $\lim_{x\to1^-}[\frac{1}{\sin(\pi x)}-\frac{1}{\pi(1-x)}]$。可令 $t=\pi(1-x)$，使用 $t-\sin t\sim\frac{t^3}{6}$。
\end{exercise}
\begin{exercise}{347}{常规·递推极限}
设 $0<x_1<\frac{\pi}{2}$，$x_{n+1}=\sin x_n$。证明 $x_n$ 单调递减且有下界，求其极限。可使用 $0<\sin t<t$（$0<t<\frac{\pi}{2}$）。
\end{exercise}
\begin{exercise}{348}{基础·无穷远极限变式}
求 $\lim_{x\to+\infty}\frac{x+\sin x}{x}$。提示：使用有界因子估计。
\end{exercise}
\begin{exercise}{349}{常规·对数化极限综合}
求 $\lim_{x\to0}(1+x+2x^2)^{\frac{1}{x}}$，并求其对数的极限。
\end{exercise}
\subsection{泰勒公式}
\begin{exercise}{350}{基础·麦克劳林公式}
写出 $e^x$ 在 $0$ 处的三阶麦克劳林公式，余项用 $o(x^3)$ 表示，并用多项式估算 $e^{0.1}$。
\end{exercise}
\begin{exercise}{351}{常规·泰勒公式求极限}
用泰勒公式求 $\lim_{x\to0}\frac{x\ln(1+x)}{e^x-x-1}$。
\end{exercise}
\begin{exercise}{352}{常规·三阶相消}
用 $\sin x$、$\cos x$ 的三阶以内展开求 $\lim_{x\to0}\frac{\sin x-x\cos x}{x^3}$。
\end{exercise}
\begin{exercise}{353}{常规·对称根式展开}
求 $\lim_{x\to0}\frac{\sqrt{1+x}+\sqrt{1-x}-2}{x^2}$。提示：将两个根式展开到二阶。
\end{exercise}
\begin{exercise}{354}{常规·泰勒系数反求参数}
求常数 $a,b$，使 $\lim_{x\to0}\frac{\ln(1+x)-ax-bx^2}{x^2}=2$。
\end{exercise}
\begin{exercise}{355}{基础·根式二阶展开}
写出 $\sqrt{1+x^2}$ 在 $0$ 处到 $x^2$ 项的展开，并求 $\lim_{x\to0}\frac{\sqrt{1+x^2}-1}{x^2}$。
\end{exercise}
\begin{exercise}{356}{常规·指数与余弦比较}
求 $\lim_{x\to0}\frac{\cos x-e^{\frac{x^2}{2}}}{x^2}$。
\end{exercise}
\begin{exercise}{357}{常规·已知余项求极限}
已知 $\sin6x+xf(x)=o(x^3)$（$x\to0$）。利用 $\sin6x=6x-36x^3+o(x^3)$，求 $\lim_{x\to0}\frac{6+f(x)}{x^2}$。
\end{exercise}
\begin{exercise}{358}{基础·展开式匹配}
求 $a,b$，使 $e^x-(1+bx+ax^2)=o(x^2)$（$x\to0$）。
\end{exercise}
\begin{exercise}{359}{常规·根式差主项}
求 $\sqrt{1-2x}-\sqrt[3]{1-3x}$ 在 $x\to0$ 时的二阶主项。可使用 $(1+u)^\alpha=1+\alpha u+\frac{\alpha(\alpha-1)u^2}{2}+o(u^2)$。
\end{exercise}
\begin{exercise}{360}{基础·对数展开}
设 $f(x)=\ln(1-3x)$。写出 $f$ 在 $0$ 处的三阶麦克劳林公式，并求 $f'''(0)$。
\end{exercise}
\begin{exercise}{361}{常规·拉格朗日余项}
用 $e^x=1+x+\frac{e^\xi x^2}{2}$（$\xi$ 介于 $0$ 与 $x$ 之间）证明 $e^x>1+x$ 对 $x\ne0$ 成立。
\end{exercise}
\begin{exercise}{362}{常规·泰勒公式数值近似}
用 $\sin x$ 的三次麦克劳林多项式估算 $\sin0.1$，并利用四阶拉格朗日余项给出误差上界。
\end{exercise}
\begin{exercise}{363}{常规·二阶泰勒余项}
已知 $f\in C^2[0,1]$，$f(0)=f'(0)=0$、$|f''(x)|\le2$。用泰勒公式证明 $|f(x)|\le x^2$（$0\le x\le1$）。
\end{exercise}
\begin{exercise}{364}{基础·泰勒展开变式}
写出 $\cos2x$ 在 $0$ 处到 $x^2$ 项的展开，余项用 $o(x^2)$ 表示。
\end{exercise}
\begin{exercise}{365}{常规·展开求极限综合}
用泰勒公式求 $\lim_{x\to0}\frac{e^{2x}-1-2x}{1-\cos2x}$。
\end{exercise}
\subsection{函数的单调性与曲线的凹凸性}
\begin{exercise}{366}{基础·单调区间}
求 $f(x)=x-\frac32x^{\frac{2}{3}}$ 的单调区间，其中分数幂按实立方根定义。提示：分别在 $x<0$ 与 $x>0$ 上求导，并检查 $x=0$。
\end{exercise}
\begin{exercise}{367}{常规·导数证明不等式}
用导数证明 $\tan x>x$（$0<x<\frac{\pi}{2}$）。提示：考察 $F(x)=\tan x-x$。
\end{exercise}
\begin{exercise}{368}{常规·对数函数比较}
用导数证明 $\ln(1+x)>\frac{x}{1+x}$（$x>0$）。提示：考察两边之差。
\end{exercise}
\begin{exercise}{369}{基础·单调性与导数}
求 $f(x)=x^3$ 的导数，说明 $f$ 在实轴上的单调性，并计算 $f'(0)$。
\end{exercise}
\begin{exercise}{370}{常规·方程唯一根}
证明 $x+\sqrt x-1=0$ 在 $(0,1)$ 内有唯一实根。提示：使用零点定理和单调性。
\end{exercise}
\begin{exercise}{371}{常规·三次方程根数}
利用 $f(x)=x^3-3x$ 的单调区间，求方程 $x^3-3x=0$ 的实根，并判断方程 $x^3-3x=3$ 的实根个数。
\end{exercise}
\begin{exercise}{372}{常规·高次方程唯一根}
用导数与无穷远极限证明方程 $x^5+x+1=0$ 有唯一实根。
\end{exercise}
\begin{exercise}{373}{常规·商函数单调性}
求 $F(x)=\frac{\ln(1+x)}{x}$（$x>0$）的单调性。可使用 $\ln(1+x)>\frac{x}{1+x}$。
\end{exercise}
\begin{exercise}{374}{常规·指数函数比较}
用导数证明 $(e+x)^e<e^{e+x}$（$x>0$）。提示：取对数后考察 $e+x-e\ln(e+x)$。
\end{exercise}
\begin{exercise}{375}{常规·中值定理与导数下界}
设 $1<a<b<2$。用中值定理证明 $\ln b-\ln a>\frac{b-a}{2}$。
\end{exercise}
\begin{exercise}{376}{基础·凹凸性与拐点}
求 $f(x)=x^3-3x^2$ 的凹凸区间和拐点，注明各区间内 $f''$ 的符号。
\end{exercise}
\begin{exercise}{377}{基础·三角函数凹凸性}
求 $f(x)=1+\sin x$ 在 $(\pi,\frac{3\pi}{2})$ 上的单调性与凹凸性，注明 $f'$、$f''$ 的符号。
\end{exercise}
\begin{exercise}{378}{常规·指数中点不等式}
利用 $(e^{\frac{x}{2}}-e^{\frac{y}{2}})^2\ge0$ 证明 $\frac{e^x+e^y}{2}\ge e^{\frac{x+y}{2}}$，并求等号条件。
\end{exercise}
\begin{exercise}{379}{常规·分段曲线凹凸性}
将 $f(x)=|x^2-1|$ 写为分段函数，求各段的二阶导数和凹凸性，并判断 $x=\pm1$ 处是否可导。
\end{exercise}
\begin{exercise}{380}{基础·拐点条件}
求曲线 $y=2x^3-3x^2+1$ 的拐点坐标，并验证二阶导数在其两侧变号。
\end{exercise}
\begin{exercise}{381}{常规·偶函数导数符号}
设 $f\in C^2(\mathbb R)$ 为偶函数，在 $(-\infty,0)$ 上 $f'>0$、$f''<0$。使用 $f'$ 为奇函数、$f''$ 为偶函数，求它们在 $(0,+\infty)$ 上的符号。
\end{exercise}
\begin{exercise}{382}{基础·水平拐点}
求 $f(x)=x^3$ 的驻点与拐点，并判断驻点是否为极值点。
\end{exercise}
\begin{exercise}{383}{常规·导数关系与凹凸性}
设 $f$ 二阶可导，$f''(x)+[f'(x)]^2=x$ 且 $f'(0)=0$。求 $f''(0)$。若另知 $f'''(0)=1$ 且 $f'''$ 在 $0$ 附近连续，判断 $(0,f(0))$ 是否为拐点。
\end{exercise}
\begin{exercise}{384}{基础·函数与导函数}
设 $f(x)=x^3-3x$。求 $f'$ 的零点及符号，并据此写出 $f$ 的单调区间。
\end{exercise}
\begin{exercise}{385}{基础·单调性变式}
求 $f(x)=x^3-6x$ 的单调区间和驻点。
\end{exercise}
\begin{exercise}{386}{常规·凹凸与零点综合}
设 $f\in C^2[0,1]$，$f(0)=f(1)=0$，且 $f''>0$ 在 $(0,1)$ 成立。先用罗尔定理证明 $f'$ 有零点，再利用 $f'$ 严格递增说明该零点唯一。
\end{exercise}
\subsection{函数的极值和最值}
\begin{exercise}{387}{基础·二阶极值判别}
求 $f(x)=x^3-3x^2+1$ 的驻点，并用二阶导数判别极值。
\end{exercise}
\begin{exercise}{388}{常规·不可导点与极值}
求 $f(x)=2x-\frac32x^{\frac{2}{3}}+1$ 的极值，其中分数幂按实立方根定义。提示：同时检查 $f'=0$ 的点和 $x=0$。
\end{exercise}
\begin{exercise}{389}{基础·指数乘积极值}
求 $f(x)=x2^x$ 的极小值点和极小值。
\end{exercise}
\begin{exercise}{390}{常规·由导函数判别极值}
设 $f'(x)=x(x+2)(x-1)^2(x-3)$。求 $f$ 的极大值点、极小值点，并判断 $x=1$ 是否为极值点。
\end{exercise}
\begin{exercise}{391}{基础·驻点分类}
已知 $f$ 二阶可导，$f'(2)=0$、$f''(2)=3$。判断 $x=2$ 处的极值类型，说明使用的判别法。
\end{exercise}
\begin{exercise}{392}{常规·导数极限与极值}
设 $f$ 二阶可导且 $f'(a)=0$，$\lim_{x\to a}\frac{f'(x)}{x-a}=-1$。求 $f''(a)$，判断 $a$ 处极值类型。
\end{exercise}
\begin{exercise}{393}{常规·增量符号与极值}
设 $\lim_{x\to a}\frac{f(x)-f(a)}{(x-a)^2}=-1$。由极限保号性判断 $a$ 处的极值类型，并求 $f'(a)$。
\end{exercise}
\begin{exercise}{394}{常规·分段函数极值}
设 $f(x)=x^2+1$（$x>0$），$f(x)=xe^x+1$（$x\le0$）。求全部极值点和极值。
\end{exercise}
\begin{exercise}{395}{常规·隐函数极值}
在方程 $x^2+y^2=1$ 的上半圆分支 $y>0$ 上，求 $y$ 关于 $x$ 的极值及取值点。
\end{exercise}
\begin{exercise}{396}{常规·参数与极值}
求 $a$ 的范围，使 $f(x)=x^3+ax$ 在实轴上没有局部极值。提示：考察 $f'(x)=3x^2+a$。
\end{exercise}
\begin{exercise}{397}{基础·闭区间最值}
求 $f(x)=x+\sqrt{1-x}$ 在 $[-5,1]$ 上的最大值和最小值。
\end{exercise}
\begin{exercise}{398}{常规·圆柱体积最值}
将半径为 $1$ 的球内接正圆柱的高记为 $h$（$0<h<2$）。利用 $r^2+(\frac{h}{2})^2=1$ 建立体积函数，求最大体积及对应的高。
\end{exercise}
\begin{exercise}{399}{常规·利润最值}
设总成本 $C(Q)=6Q+200$，总收益 $R(Q)=12Q-0.01Q^2$，$Q\ge0$ 为连续产量。求利润最大时的产量和最大利润。
\end{exercise}
\begin{exercise}{400}{常规·收入成本与最值}
设 $R(x)=26x-2x^2$、$C(x)=8x+x^2$，$x\ge0$。求利润函数及最大利润对应的产量。
\end{exercise}
\begin{exercise}{401}{常规·最值证明不等式}
利用 $F(x)=\frac{x^2}{2}+\frac{1}{2}-x$ 的最小值证明 $\frac{x^2}{2}+\frac{1}{2}\ge x$，并求等号条件。
\end{exercise}
\begin{exercise}{402}{常规·凸性与最小值}
设 $f\in C^2(\mathbb R)$，$f(0)=0$、$f'(0)=1$ 且 $f''>0$。考察 $F(x)=f(x)-x$ 的导数，证明 $f(x)\ge x$，求等号条件。
\end{exercise}
\begin{exercise}{403}{基础·闭区间最值变式}
求 $f(x)=x^2-x$ 在 $[0,1]$ 上的最大值与最小值及取值点。
\end{exercise}
\begin{exercise}{404}{常规·无盖箱体最值}
从边长为 $6$ 的正方形薄板四角各剪去边长为 $x$ 的小正方形，折成无盖长方体。写出体积函数与可行域，求最大体积及对应的 $x$。
\end{exercise}
\subsection{函数图形的描绘}
\begin{exercise}{405}{基础·水平与铅直渐近线}
求 $y=\frac{x+1}{x}$ 的水平渐近线和铅直渐近线。
\end{exercise}
\begin{exercise}{406}{基础·有理函数渐近线}
求 $y=\frac{1}{1-x^2}$ 的水平渐近线和铅直渐近线。
\end{exercise}
\begin{exercise}{407}{常规·可去间断与渐近线}
求 $f(x)=\frac{x^2-x-2}{x^2-1}$ 的水平、铅直渐近线，并说明 $x=-1$ 是否为铅直渐近线。
\end{exercise}
\begin{exercise}{408}{基础·斜渐近线}
用多项式除法求 $y=\frac{x^2}{x-1}$ 的斜渐近线和铅直渐近线。
\end{exercise}
\begin{exercise}{409}{常规·指数复合渐近线}
求 $y=(2x-1)e^{\frac{1}{x}}$ 的斜渐近线及 $x\to0^+$ 的极限。可使用 $e^{\frac{1}{x}}=1+\frac{1}{x}+o(\frac{1}{x})$（$|x|\to\infty$）。
\end{exercise}
\begin{exercise}{410}{基础·振荡因子与渐近线}
设 $f(x)=x\sin(\frac{1}{x})$（$x>0$）。求 $x\to+\infty$ 时的水平渐近线，并求 $x\to0^+$ 的极限。
\end{exercise}
\begin{exercise}{411}{常规·指数曲线渐近线}
求 $f(x)=xe^{\frac{1}{x}}$ 的斜渐近线及铅直渐近线。提示：分别计算 $x\to0^\pm$ 的极限。
\end{exercise}
\begin{exercise}{412}{常规·对数函数水平渐近线}
分别求 $f(x)=\frac{\ln(1+e^x)}{x}$ 在正、负无穷方向的水平渐近线。
\end{exercise}
\begin{exercise}{413}{基础·反三角函数渐近线}
求 $f(x)=\arctan x$ 在正、负无穷方向的水平渐近线，并求它与坐标轴的交点。
\end{exercise}
\begin{exercise}{414}{常规·有理函数图形分析}
对 $f(x)=\frac{x^2}{x-1}$，求定义域、渐近线、单调区间和极值，据此画出图形草图。
\end{exercise}
\begin{exercise}{415}{常规·周期函数图形分析}
求 $f(x)=\cos x$ 的最小正周期，在 $[-\pi,\pi]$ 上求单调区间和极值，并画出该区间的图形草图。
\end{exercise}
\begin{exercise}{416}{基础·斜渐近线变式}
求 $f(x)=x+\frac{\sin x}{x}$（$x\ne0$）的斜渐近线，并求 $x\to0$ 时的极限。
\end{exercise}
\begin{exercise}{417}{常规·渐近线与单调性综合}
求 $f(x)=\frac{x^2+1}{x-1}$ 的铅直、斜渐近线及单调区间。
\end{exercise}
\subsection{曲率}
\begin{exercise}{418}{基础·显函数曲率}
用 $K=\frac{|y''|}{(1+(y')^2)^{\frac{3}{2}}}$ 求曲线 $y=\frac{1}{x}$ 在 $(1,1)$ 处的曲率与曲率半径。
\end{exercise}
\begin{exercise}{419}{常规·参数曲线曲率}
求曲线 $x=2\cos^3t$、$y=2\sin^3t$ 在 $t=\frac{\pi}{4}$ 处的曲率半径。可使用 $K=\frac{|x'y''-y'x''|}{((x')^2+(y')^2)^{\frac{3}{2}}}$。
\end{exercise}
\begin{exercise}{420}{常规·曲率与曲率圆}
求抛物线 $y=x^2$ 在原点处的曲率和曲率半径，并写出曲率圆方程。提示：圆心在 $y$ 轴正半轴上。
\end{exercise}
\begin{exercise}{421}{基础·根式曲线曲率}
求 $y=\sqrt x$ 在 $(1,1)$ 处的曲率半径。
\end{exercise}
\begin{exercise}{422}{基础·抛物线曲率变式}
求 $y=2x^2$ 在原点处的曲率及曲率半径。
\end{exercise}
\begin{exercise}{423}{常规·椭圆曲率综合}
设椭圆 $x=2\cos t$、$y=\sin t$。求 $t=0$ 与 $t=\frac{\pi}{2}$ 两点的曲率半径并比较大小。可使用参数曲线曲率公式。
\end{exercise}

\clearpage
\section{第4章 不定积分}
\subsection{不定积分的概念与性质}
\begin{exercise}{424}{基础·原函数关系}
设 $f$ 在 $(1,+\infty)$ 上的一个原函数为 $F(x)=\ln(x+\sqrt{x^2-1})$。求 $f(x)$。
\end{exercise}
\begin{exercise}{425}{基础·原函数与常数}
已知 $\ln x$ 是 $f$ 在 $(0,+\infty)$ 上的一个原函数。求 $f(x)$，并判断 $\ln(2x)$ 是否也是它的原函数。
\end{exercise}
\begin{exercise}{426}{常规·两个原函数之差}
在 $(0,1)$ 上求 $F(x)=\arcsin(2x-1)$、$G(x)=2\arcsin\sqrt x$ 的导数，说明它们为何相差常数，并求这个常数。
\end{exercise}
\begin{exercise}{427}{常规·分段原函数}
设 $f(x)=3x^2$（$x>0$），$f(x)=2x$（$x\le0$）。验证 $F(x)=x^3$（$x>0$）、$F(x)=x^2$（$x\le0$）是 $f$ 在实轴上的原函数，单独检查 $x=0$。
\end{exercise}
\begin{exercise}{428}{基础·连续函数的原函数}
说明 $f(x)=|x|$ 在实轴上存在原函数的依据，并验证 $F(x)=x|x|/2$ 是它的一个原函数。
\end{exercise}
\begin{exercise}{429}{常规·分段原函数拼接}
求 $f(x)=x$（$x<0$）、$f(x)=2x$（$x\ge0$）在实轴上满足 $F(0)=0$ 的原函数。
\end{exercise}
\begin{exercise}{430}{基础·原函数与导函数}
已知 $\sin x$ 是 $f(x)$ 的一个原函数。求 $f(x)$ 及 $\int f(x)\,\dd x$。
\end{exercise}
\begin{exercise}{431}{常规·斜率与初值}
一条曲线经过 $(e,-1)$，在每个 $x>0$ 处的切线斜率为 $\ln x/x$。求曲线方程。
\end{exercise}
\begin{exercise}{432}{基础·积分与微分}
设 $f$ 可导。分别求 $\int f'(x)\,\dd x$ 和 $\dd[\int f'(x)\,\dd x]$。
\end{exercise}
\begin{exercise}{433}{基础·积分与求导}
分别计算 $\int3x^2\,\dd x$、$\dd(\int3x^2\,\dd x)$ 及 $(\int3x^2\,\dd x)'$。
\end{exercise}
\begin{exercise}{434}{常规·原函数之差}
设 $F'=G'$ 在开区间 $I$ 上成立。用中值定理说明 $F-G$ 为常数；若 $F(x_0)=G(x_0)$，求这个常数。
\end{exercise}
\begin{exercise}{435}{基础·微分的积分}
求 $\int\dd(x^2+\sin x)$，并对结果求导核验。
\end{exercise}
\begin{exercise}{436}{常规·换元恢复导函数}
已知 $f'(\sin^2t)=\cos2t$（$0<t<\pi/2$），且 $f(1/2)=0$。先求 $f'(x)$（$0<x<1$），再求 $f(x)$。
\end{exercise}
\begin{exercise}{437}{基础·拆项积分}
求 $\int(x-1/x-3/\sqrt{1-x^2})\,\dd x$，在 $(-1,0)$ 或 $(0,1)$ 上分别考虑。
\end{exercise}
\begin{exercise}{438}{基础·三角恒等变形}
求 $\int\cos2x/(\sin^2x\cos^2x)\,\dd x$，在 $\sin x\cos x\ne0$ 的区间上考虑。提示：将分子写成 $\cos^2x-\sin^2x$。
\end{exercise}
\begin{exercise}{439}{基础·三角基本积分}
求 $\int(1-2\cos x)/\sin^2x\,\dd x$，在不含 $k\pi$ 的开区间上考虑。
\end{exercise}
\begin{exercise}{440}{基础·指数化简积分}
求 $\int[(5e)^{-x}+1-e^x]\,\dd x$。
\end{exercise}
\begin{exercise}{441}{基础·降幂积分}
求 $\int[6\sqrt x+5e^x+\cos^2(x/2)]\,\dd x$（$x>0$）。
\end{exercise}
\begin{exercise}{442}{常规·多项式除法积分}
求 $\int[x^4/(1+x^2)-2/\sqrt{1-x^2}]\,\dd x$（$|x|<1$）。提示：先对有理项作除法。
\end{exercise}
\begin{exercise}{443}{基础·对数原函数}
在 $(0,+\infty)$ 上求 $\int\dd x/x$，并求满足 $F(1)=2$ 的原函数。
\end{exercise}
\begin{exercise}{444}{常规·分段原函数小综合}
设 $f(x)=x$（$x<0$）、$f(x)=2x$（$x\ge0$）。求它在实轴上满足 $F(0)=1$ 的原函数，并检查 $F'(0)$。
\end{exercise}
\subsection{换元积分法}
\begin{exercise}{445}{基础·凑微分}
求 $\int\dd x/\sqrt{2x+3}$（$x>-3/2$）及 $\int(x+1)/(x^2+2x+5)\,\dd x$。
\end{exercise}
\begin{exercise}{446}{基础·根式换元}
求 $\int e^{2\sqrt x}/\sqrt x\,\dd x$ 与 $\int\sin\sqrt x/\sqrt x\,\dd x$（$x>0$）。提示：令 $t=\sqrt x$。
\end{exercise}
\begin{exercise}{447}{基础·复合函数换元}
求 $\int x^2e^{x^3}\,\dd x$ 和 $\int x/(1+x^4)\,\dd x$。
\end{exercise}
\begin{exercise}{448}{基础·三角换元}
求 $\int\sin^2x\cos^2x\,\dd x$ 与 $\int\sin^3x\cos x\,\dd x$。
\end{exercise}
\begin{exercise}{449}{基础·平方和与平方差积分}
求 $\int\dd x/(4+x^2)$ 和 $\int\dd x/(4-x^2)$，后一式在不含 $\pm2$ 的开区间上考虑。
\end{exercise}
\begin{exercise}{450}{基础·指数复合换元}
求 $\int e^{\tan x}\sec^2x\,\dd x$ 与 $\int(1-x^{-2})e^{x+1/x}\,\dd x$，分别在被积函数连续的开区间上考虑。
\end{exercise}
\begin{exercise}{451}{基础·对数换元}
求 $\int\dd x/[x\sqrt{1-\ln^2x}]$（$e^{-1}<x<e$）及 $\int(1+\ln x)/(1+x^2\ln^2x)\,\dd x$（$x>0$）。
\end{exercise}
\begin{exercise}{452}{基础·反三角函数换元}
求 $\int\arcsin x/\sqrt{1-x^2}\,\dd x$（$|x|<1$）和 $\int\arcsin\sqrt x/\sqrt{x(1-x)}\,\dd x$（$0<x<1$）。
\end{exercise}
\begin{exercise}{453}{常规·配方换元}
求 $\int\dd x/\sqrt{x-x^2}$（$0<x<1$）。提示：配方 $x-x^2=1/4-(x-1/2)^2$。
\end{exercise}
\begin{exercise}{454}{基础·三角内函数}
求 $\int\cos x/(9+\sin^2x)\,\dd x$，并求满足 $F(0)=0$ 的原函数。
\end{exercise}
\begin{exercise}{455}{基础·根式拆项积分}
求 $\int(1-x)/\sqrt{9-4x^2}\,\dd x$（$|x|<3/2$）。
\end{exercise}
\begin{exercise}{456}{常规·立方根换元}
求 $\int(\sin x+\cos x)/\sqrt[3]{\sin x-\cos x}\,\dd x$，立方根按实数定义，在 $\sin x-\cos x\ne0$ 的开区间上考虑。
\end{exercise}
\begin{exercise}{457}{常规·根式三角复合}
求 $\int x\tan(\sqrt{1-x^2})/\sqrt{1-x^2}\,\dd x$（$|x|<1$）。提示：令 $u=\sqrt{1-x^2}$。
\end{exercise}
\begin{exercise}{458}{常规·对数根式换元}
求 $\int[\ln(x+\sqrt{1+x^2})+1]/\sqrt{1+x^2}\,\dd x$。提示：令 $u=\ln(x+\sqrt{1+x^2})$。
\end{exercise}
\begin{exercise}{459}{常规·指数根式换元}
求 $\int e^{\sqrt{1+\sin x}}\cos x/\sqrt{1+\sin x}\,\dd x$，在 $1+\sin x>0$ 的开区间上考虑。提示：令 $u=\sqrt{1+\sin x}$。
\end{exercise}
\begin{exercise}{460}{基础·拆项与反正切换元}
求 $\int[x+(\arctan x)^2]/(1+x^2)\,\dd x$。
\end{exercise}
\begin{exercise}{461}{常规·三角幂有理化}
求 $\int\sin x\cos^3x/(1+\cos^2x)\,\dd x$。提示：令 $u=\cos x$，再作多项式除法。
\end{exercise}
\begin{exercise}{462}{常规·正切换元}
在 $(-\pi/2,\pi/2)$ 上求 $\int\dd x/(\sin^2x+2\cos^2x)$。提示：令 $t=\tan x$。
\end{exercise}
\begin{exercise}{463}{基础·指数换元}
求 $\int\dd x/(1+e^x)$，并求满足 $F(0)=0$ 的原函数。
\end{exercise}
\begin{exercise}{464}{常规·对数差整体换元}
求 $\int[\ln(x+1)-\ln x]/[x(x+1)]\,\dd x$（$x>0$）。提示：令 $u=\ln(x+1)-\ln x$。
\end{exercise}
\begin{exercise}{465}{基础·对数换元与拆分}
求 $\int\ln^2x/[x(1+\ln^2x)]\,\dd x$（$x>0$）。提示：令 $t=\ln x$。
\end{exercise}
\begin{exercise}{466}{常规·根式有理化}
求 $\int\dd x/(x-\sqrt{x^2-1})$（$x>1$）。提示：用共轭式将被积函数化为 $x+\sqrt{x^2-1}$。
\end{exercise}
\begin{exercise}{467}{基础·指数内函数凑微分}
求 $\int e^{\sin2x}\cos2x\,\dd x$。
\end{exercise}
\begin{exercise}{468}{常规·对数正切换元}
在 $0<x<\pi/2$ 上求 $\int\ln(\tan x)/\sin2x\,\dd x$。提示：计算 $\ln\tan x$ 的导数。
\end{exercise}
\begin{exercise}{469}{常规·商函数导数识别}
求 $\int(x\cos x-\sin x)/x^2\,\dd x$，在两个半轴上分别考虑。提示：计算 $(\sin x/x)'$。
\end{exercise}
\begin{exercise}{470}{常规·三角指数整体换元}
求 $\int e^{\sin x}\cos x/(1+e^{\sin x})\,\dd x$。提示：令 $u=e^{\sin x}$。
\end{exercise}
\begin{exercise}{471}{基础·半角积分}
求 $\int\sin x/(1+\cos x)\,\dd x$，在 $1+\cos x\ne0$ 的开区间上考虑。
\end{exercise}
\begin{exercise}{472}{常规·已知原函数换元}
已知 $\int f(u)\,\dd u=\ln(\sin u)+C$（$0<u<\pi$）。求 $\int xf(1-x^2)\,\dd x$，在 $|x|<1$ 上考虑。
\end{exercise}
\begin{exercise}{473}{常规·抽象函数商的积分}
设 $f\in C^1(I)$ 且 $0\notin I$。利用 $[f(x)/x]'$ 求 $\int[xf'(x)-f(x)]/x^2\,\dd x$。
\end{exercise}
\begin{exercise}{474}{常规·根式换元综合}
求 $\int x^3\sqrt{4-x^2}\,\dd x$（$|x|<2$）。提示：令 $u=4-x^2$，将余下的 $x^2$ 写成 $4-u$。
\end{exercise}
\begin{exercise}{475}{常规·平方和根式积分}
求 $\int x^2/(1+x^2)^{3/2}\,\dd x$。提示：拆成 $1/\sqrt{1+x^2}-1/(1+x^2)^{3/2}$。
\end{exercise}
\begin{exercise}{476}{常规·平方差根式积分}
求 $\int x/(x^2-4)^{3/2}\,\dd x$（$x>2$）。提示：令 $u=x^2-4$。
\end{exercise}
\begin{exercise}{477}{常规·立方根有理化}
求 $\int\dd x/[1+\sqrt[3]{x+1}]$（$x>-1$）。提示：令 $t=\sqrt[3]{x+1}$，将 $3t^2/(1+t)$ 作多项式除法。
\end{exercise}
\begin{exercise}{478}{常规·不同根指数换元}
求 $\int\dd x/(\sqrt x+\sqrt[3]{x^2})$（$x>0$）。提示：令 $t=x^{1/6}$，将积分化为 $\int6t^2/(1+t)\,\dd t$。
\end{exercise}
\begin{exercise}{479}{常规·根式分式换元}
求 $\int\dd x/[(x+1)\sqrt x]$（$x>0$）。提示：令 $t=\sqrt x$。
\end{exercise}
\begin{exercise}{480}{基础·指数有理积分}
求 $\int\dd x/(e^x+e^{2x})$。提示：令 $t=e^x$。
\end{exercise}
\begin{exercise}{481}{常规·指数根式换元}
求 $\int\dd x/\sqrt{e^x+1}$。提示：令 $t=\sqrt{e^x+1}$，则 $\dd x=2t\,\dd t/(t^2-1)$。
\end{exercise}
\begin{exercise}{482}{常规·幂换元与拆分}
求 $\int\dd x/[x(x^2+1)]$，分别在两个半轴上考虑。提示：拆成 $1/x-x/(x^2+1)$。
\end{exercise}
\begin{exercise}{483}{基础·倒数换元}
求 $\int\sin(1/x)/x^2\,\dd x$（$x\ne0$）。
\end{exercise}
\begin{exercise}{484}{常规·混合内函数换元}
求 $\int(1+x)/(x^2+x\ln x)\,\dd x$，在 $x>1$ 上考虑。提示：令 $u=x+\ln x$。
\end{exercise}
\begin{exercise}{485}{常规·配方与三角代换}
求 $\int\sqrt{2x-x^2}\,\dd x$（$0<x<2$）。提示：先配方为 $1-(x-1)^2$。
\end{exercise}
\begin{exercise}{486}{常规·指数根式复合}
求 $\int e^x/\sqrt{1+e^{2x}}\,\dd x$。提示：令 $t=e^x$。
\end{exercise}
\begin{exercise}{487}{常规·三角代换}
求 $\int x^2/\sqrt{1-x^2}\,\dd x$（$|x|<1$）。提示：令 $x=\sin t$，$|t|<\pi/2$，再降幂。
\end{exercise}
\begin{exercise}{488}{基础·换元积分变式}
求 $\int x/(1+3x^2)\,\dd x$。
\end{exercise}
\begin{exercise}{489}{常规·换元与初值综合}
求 $\int x/\sqrt{1+x^2}\,\dd x$，并求满足 $F(0)=2$ 的原函数。
\end{exercise}
\subsection{分部积分法}
\begin{exercise}{490}{基础·多项式乘指数或三角函数}
求 $\int xe^{2x}\,\dd x$ 和 $\int x^2\cos x\,\dd x$。
\end{exercise}
\begin{exercise}{491}{基础·反三角与对数分部积分}
求 $\int\ln x\,\dd x$（$x>0$）和 $\int\arctan x\,\dd x$。
\end{exercise}
\begin{exercise}{492}{常规·循环积分}
用两次分部积分求 $\int e^{-x}\cos x\,\dd x$。
\end{exercise}
\begin{exercise}{493}{基础·移位对数积分}
求 $\int x\ln(x+1)\,\dd x$（$x>-1$）。
\end{exercise}
\begin{exercise}{494}{基础·反正弦分部积分}
求 $\int\arcsin(x+1)\,\dd x$（$-2<x<0$）。
\end{exercise}
\begin{exercise}{495}{常规·凑微分与分部配合}
求 $\int x\cot x/\sin^2x\,\dd x$，在 $\sin x\ne0$ 的开区间上考虑。提示：$\cot x/\sin^2x=-(\cot^2x)'/2$。
\end{exercise}
\begin{exercise}{496}{基础·三角函数分部积分}
求 $\int x/\cos^2x\,\dd x$，在 $\cos x\ne0$ 的开区间上考虑。
\end{exercise}
\begin{exercise}{497}{常规·对数三角分部积分}
在 $0<x<\pi$ 上求 $\int\ln(\sin x)/\sin^2x\,\dd x$。提示：取 $u=\ln\sin x$、$\dd v=\csc^2x\,\dd x$。
\end{exercise}
\begin{exercise}{498}{常规·根式分部积分}
用分部积分求 $\int x^3/\sqrt{1+x^2}\,\dd x$。提示：取 $u=x^2$、$\dd v=x\,\dd x/\sqrt{1+x^2}$。
\end{exercise}
\begin{exercise}{499}{常规·指数三角循环积分}
用两次分部积分求 $\int e^{2x}\sin3x\,\dd x$。
\end{exercise}
\begin{exercise}{500}{常规·平方和根式积分}
求 $\int\sqrt{x^2+1}\,\dd x$。提示：分部后使用 $x^2=(x^2+1)-1$，将同类积分移项。
\end{exercise}
\begin{exercise}{501}{常规·平方差根式积分}
求 $\int\sqrt{4-x^2}\,\dd x$（$|x|<2$）。可令 $x=2\sin t$。
\end{exercise}
\begin{exercise}{502}{常规·反正弦乘多项式}
求 $\int x\arcsin x\,\dd x$（$|x|<1$）。提示：先取 $u=\arcsin x$，再对余下的根式积分作三角代换。
\end{exercise}
\begin{exercise}{503}{常规·反正切乘负幂}
求 $\int\arctan x/x^2\,\dd x$，在两个半轴上分别考虑。可使用 $1/[x(1+x^2)]=1/x-x/(1+x^2)$。
\end{exercise}
\begin{exercise}{504}{基础·对数乘有理因子}
求 $\int\ln x/x^2\,\dd x$（$x>0$）。
\end{exercise}
\begin{exercise}{505}{常规·指数换元与分部积分}
求 $\int xe^{\sqrt x}/\sqrt x\,\dd x$（$x>0$）。提示：令 $t=\sqrt x$，积分化为 $2\int t^2e^t\,\dd t$。
\end{exercise}
\begin{exercise}{506}{常规·对数与根式分部积分}
求 $\int\ln x/\sqrt x\,\dd x$（$x>0$）。提示：取 $u=\ln x$、$\dd v=x^{-1/2}\,\dd x$。
\end{exercise}
\begin{exercise}{507}{常规·反正切根式积分}
求 $\int\arctan\sqrt x\,\dd x$（$x>0$）。提示：分部积分后令 $t=\sqrt x$。
\end{exercise}
\begin{exercise}{508}{常规·指数有理导数识别}
求 $\int e^x(x-1)/x^2\,\dd x$（$x\ne0$）。提示：计算 $(e^x/x)'$。
\end{exercise}
\begin{exercise}{509}{常规·三角商的逆向求导}
求 $\int[(1+x)\sin x+\cos x]/(1+x)^2\,\dd x$（$x\ne-1$）。提示：计算 $[\cos x/(1+x)]'$。
\end{exercise}
\begin{exercise}{510}{常规·恢复导函数与分部积分}
已知 $f'(e^t)=\cos t$ 对所有实数 $t$ 成立且 $f(1)=0$。求 $f(x)$（$x>0$）。提示：先恢复 $f'(x)$，再令 $u=\ln x$。
\end{exercise}
\begin{exercise}{511}{常规·正割三次积分}
在 $(-\pi/2,\pi/2)$ 上用分部积分求 $\int\sec^3x\,\dd x$。提示：取 $u=\sec x$、$\dd v=\sec^2x\,\dd x$，再使用 $\tan^2x=\sec^2x-1$。
\end{exercise}
\begin{exercise}{512}{基础·分部积分变式}
求 $\int x\sin x\,\dd x$。
\end{exercise}
\begin{exercise}{513}{常规·对数平方积分}
求 $\int(\ln x)^2\,\dd x$（$x>0$），并确定满足 $F(1)=0$ 的原函数。
\end{exercise}
\subsection{有理函数的积分}
\begin{exercise}{514}{基础·假分式积分}
求 $\int x^3/(x^2+1)\,\dd x$。提示：先作多项式除法。
\end{exercise}
\begin{exercise}{515}{基础·两个一次因子}
求 $\int(3x-2)/(x^2+3x-4)\,\dd x$，在不含 $-4,1$ 的开区间上考虑。
\end{exercise}
\begin{exercise}{516}{常规·三个一次因子}
求 $\int\dd x/[(x+1)(x-1)(x+2)]$，在分母非零的开区间上考虑。
\end{exercise}
\begin{exercise}{517}{基础·简单部分分式}
求 $\int\dd x/[x(x-1)]$，在不含 $0,1$ 的开区间上考虑。
\end{exercise}
\begin{exercise}{518}{常规·重一次因子}
求 $\int(x^2+1)/[(x+1)^2(x-1)]\,\dd x$。提示：设分解为 $A/(x-1)+B/(x+1)+C/(x+1)^2$。
\end{exercise}
\begin{exercise}{519}{常规·一次与二次因子}
求 $\int\dd x/[(x^2+1)(1+2x)]$。提示：分解为 $A/(1+2x)+(Bx+C)/(x^2+1)$。
\end{exercise}
\begin{exercise}{520}{基础·配方与分子拆分}
求 $\int(3x+1)/(x^2+2x+5)\,\dd x$。
\end{exercise}
\begin{exercise}{521}{常规·重二次分母}
求 $\int\dd x/(x^2+1)^2$。提示：令 $x=\tan t$。
\end{exercise}
\begin{exercise}{522}{常规·四次分母分解}
求 $\int(2x^3+2x^2+5x+5)/(x^4+5x^2+4)\,\dd x$。提示：分母为 $(x^2+1)(x^2+4)$，可先把分子拆为奇、偶两部分。
\end{exercise}
\begin{exercise}{523}{常规·平移后重二次积分}
求 $\int(x+1)/(x^2+2x+2)^2\,\dd x$。提示：令 $u=x^2+2x+2$。
\end{exercise}
\begin{exercise}{524}{基础·高次式换元}
求 $\int x^2/(1+x^6)\,\dd x$。提示：令 $u=x^3$。
\end{exercise}
\begin{exercise}{525}{常规·重零点分解}
求 $\int\dd x/[x^2(1-x^2)]$，在不含 $-1,0,1$ 的开区间上考虑。提示：先拆为 $1/x^2+1/(1-x^2)$。
\end{exercise}
\begin{exercise}{526}{基础·重因子换元}
求 $\int x^3/(1+x^4)^2\,\dd x$。
\end{exercise}
\begin{exercise}{527}{常规·四次分母典型模型}
求 $\int(x^2+1)/(x^4+3x^2+2)\,\dd x$。提示：先分解分母并约分。
\end{exercise}
\begin{exercise}{528}{基础·原函数奇偶性实例}
求 $f(x)=x^3$ 满足 $F(0)=0$ 的原函数，并判断 $F$ 的奇偶性。
\end{exercise}
\begin{exercise}{529}{基础·连续函数原函数}
设 $f(x)=x^2+1$。求满足 $F(0)=0$ 的原函数，并说明 $f$ 在 $[-1,1]$ 上有最大、最小值的依据。
\end{exercise}
\begin{exercise}{530}{常规·分段原函数参数}
设 $f(x)=2(x-1)$（$x<1$）、$f(x)=\ln x$（$x\ge1$）。求 $a,b$，使 $F(x)=(x-1)^2$（$x<1$）、$F(x)=x(\ln x+a)+b$（$x\ge1$）为 $f$ 在实轴上的原函数。
\end{exercise}
\begin{exercise}{531}{常规·反正切指数分部}
求 $\int e^x\arctan(e^x)\,\dd x$。提示：先令 $u=e^x$，再分部积分。
\end{exercise}
\begin{exercise}{532}{基础·反正切整体换元}
求 $\int\arctan x/(1+x^2)\,\dd x$。
\end{exercise}
\begin{exercise}{533}{常规·指数有理对数导数}
求 $\int(1+x)/[x(1+xe^x)]\,\dd x$（$x>0$）。提示：将被积函数写成 $1+1/x-(xe^x)'/(1+xe^x)$。
\end{exercise}
\begin{exercise}{534}{常规·对数商逆向求导}
求 $\int(1-\ln x)/(x-\ln x)^2\,\dd x$（$x>0$）。可使用 $x-\ln x\ge1$，并计算 $[x/(x-\ln x)]'$。
\end{exercise}
\begin{exercise}{535}{常规·根式反正切换元}
求 $\int e^x\arctan\sqrt{e^x-1}/\sqrt{e^x-1}\,\dd x$（$x>0$）。提示：令 $t=\sqrt{e^x-1}$，积分化为 $2\int\arctan t\,\dd t$。
\end{exercise}
\begin{exercise}{536}{常规·商函数导数结构}
求 $\int[x+\sin x\cos x]/(x\sin x+\cos x)^2\,\dd x$，在分母非零的开区间上考虑。提示：计算 $[\cos x/(x\sin x+\cos x)]'$。
\end{exercise}
\begin{exercise}{537}{基础·三角有理换元}
求 $\int\cos x/[\sin x(1+\sin x)]\,\dd x$，在 $\sin x\ne0,-1$ 的开区间上考虑。提示：令 $u=\sin x$。
\end{exercise}
\begin{exercise}{538}{常规·万能代换}
在 $(-\pi/2,\pi/2)$ 上求 $\int\dd x/(1+\sin x+\cos x)$。提示：令 $t=\tan(x/2)$。
\end{exercise}
\begin{exercise}{539}{常规·万能代换与反正切}
在 $(-\pi,\pi)$ 上求 $\int\dd x/(2+\cos x)$。提示：令 $t=\tan(x/2)$。
\end{exercise}
\begin{exercise}{540}{基础·分母及其导数拆分}
求 $\int(3\cos x-\sin x)/(\cos x+\sin x)\,\dd x$，在分母非零的开区间上考虑。
\end{exercise}
\begin{exercise}{541}{基础·对称拆分三角积分}
求 $\int\sin x/(\sin x+\cos x)\,\dd x$，在分母非零的开区间上考虑。
\end{exercise}
\begin{exercise}{542}{常规·根式三角代换}
求 $\int\dd x/[2x+\sqrt{1-x^2}]$，在 $0<x<1$ 上考虑。提示：令 $x=\sin t$（$0<t<\pi/2$），把所得分子写成分母及其导数的线性组合。
\end{exercise}
\begin{exercise}{543}{常规·拆项与重二次分母}
求 $\int(x^3+1)/(1+x^2)^2\,\dd x$。提示：分开处理 $x^3/(1+x^2)^2$ 与 $1/(1+x^2)^2$。
\end{exercise}
\begin{exercise}{544}{基础·反正切指数换元}
求 $\int e^{\arctan x}/(1+x^2)\,\dd x$。
\end{exercise}
\begin{exercise}{545}{常规·指数有理换元}
求 $\int\dd x/(1+e^{x/2})$。提示：令 $t=e^{x/2}$，注意 $\dd x=2\dd t/t$。
\end{exercise}
\begin{exercise}{546}{常规·三角因子有理换元}
在 $0<x<\pi$ 上求 $\int\sin x/(2+\cos x)\,\dd x$。
\end{exercise}
\begin{exercise}{547}{基础·反正弦整体换元}
求 $\int\arcsin x/\sqrt{1-x^2}\,\dd x$（$|x|<1$）。
\end{exercise}
\begin{exercise}{548}{常规·已知函数关系的积分}
设 $y=y(x)>0$ 满足 $y^2=x$（$x>0$）。求 $\int\dd x/y(x)$，并用 $x=y^2$ 换元核验。
\end{exercise}
\begin{exercise}{549}{常规·函数关系与换元}
已知 $f(\sin^2t)=t/\sin t$（$0<t<\pi/2$）。求 $\int f(x)/\sqrt{1-x}\,\dd x$（$0<x<1$）。提示：令 $x=\sin^2t$，积分化为 $\int2t\,\dd t$。
\end{exercise}
\begin{exercise}{550}{常规·已知原函数与分部积分}
已知 $\int f(x)\,\dd x=e^{x^2}+C$。求 $f(x)$ 及 $\int xf'(x)\,\dd x$。提示：分部积分可写为 $xf(x)-\int f(x)\,\dd x$。
\end{exercise}
\begin{exercise}{551}{基础·加权原函数恢复}
已知 $\int xf(x)\,\dd x=x^3+C$（$x>0$）。求 $f(x)$ 和 $\int f(x)\,\dd x$。
\end{exercise}
\begin{exercise}{552}{常规·复合导函数恢复}
已知 $f'(e^t)=te^{-t}$ 对所有实数 $t$ 成立且 $f(1)=0$。求 $f(x)$（$x>0$）。
\end{exercise}
\begin{exercise}{553}{常规·抽象函数平方导数}
设 $f\in C^1(I)$。求 $\int f(x)f'(x)\,\dd x$ 及 $\int f'(x)/(1+f(x)^2)\,\dd x$。
\end{exercise}
\begin{exercise}{554}{常规·线性三角组合平方}
求 $\int\dd x/(\sin x+\cos x)^2$，在分母非零的开区间上考虑。提示：先写成 $2\sin^2(x+\pi/4)$。
\end{exercise}
\begin{exercise}{555}{基础·部分分式变式}
求 $\int\dd x/[(x-1)(x-2)]$，在不含 $1,2$ 的开区间上考虑。
\end{exercise}
\begin{exercise}{556}{常规·有理积分小综合}
求 $\int\dd x/[x(x^2+1)^2]$，在两个半轴上分别考虑。提示：先拆为 $1/x-x/(x^2+1)-x/(x^2+1)^2$。
\end{exercise}

\clearpage
\section{第5章 定积分}
\subsection{定积分的概念与性质}
\begin{exercise}{557}{基础·积分变量}
设 $I=\int_a^b\arcsin t\dd t$，其中 $-1<a<b<1$。求 $I$ 关于上限 $b$ 和下限 $a$ 的导数。
\end{exercise}
\begin{exercise}{558}{基础·几何意义}
利用半圆面积求 $\int_{-2}^2\sqrt{4-x^2}\dd x$。
\end{exercise}
\begin{exercise}{559}{基础·黎曼和}
将 $\lim_{n\to\infty}\frac1n\sum_{k=1}^n\frac{k}{n}$ 写成定积分并求值。
\end{exercise}
\begin{exercise}{560}{基础·黎曼和}
将 $\sum_{k=1}^n\frac1{n+k}$ 写成黎曼和，求其在 $n\to\infty$ 时的极限。
\end{exercise}
\begin{exercise}{561}{基础·黎曼和}
利用定积分求 $\lim_{n\to\infty}\sum_{k=1}^n\frac{n}{n^2+k^2}$。
\end{exercise}
\begin{exercise}{562}{基础·积分比较}
不求原函数，比较 $\int_0^1x^2\dd x$ 与 $\int_0^1x^3\dd x$ 的大小，并说明依据。
\end{exercise}
\begin{exercise}{563}{基础·积分比较}
不求原函数，比较 $\int_0^1\sqrt{1+x^2}\dd x$ 与 $\int_0^1\sqrt{1+x^3}\dd x$ 的大小。
\end{exercise}
\begin{exercise}{564}{基础·积分估值}
利用 $0\le\cos^2x\le1$，给出 $\int_0^{\frac{\pi}{2}}(1+\cos^2x)\dd x$ 的上、下界。
\end{exercise}
\begin{exercise}{565}{基础·带权积分估值}
利用 $\frac{1}{\sqrt2}\le\frac{1}{\sqrt{1+x}}\le1$（$0\le x\le1$），估计 $\int_0^1\frac{x^2}{\sqrt{1+x}}\dd x$。
\end{exercise}
\begin{exercise}{566}{基础·三角积分估值}
利用 $\frac{\sqrt2}{2}\le\sin x\le1$，估计 $\int_{\frac{\pi}{4}}^{\frac{\pi}{2}}\sin x\dd x$。
\end{exercise}
\begin{exercise}{567}{常规·积分中值定理}
设 $f$ 在 $[0,1]$ 上连续。用积分中值定理说明存在 $\xi\in[0,1]$ 使 $\int_0^1f(x)\dd x=f(\xi)$；当 $f(x)=x^2$ 时求 $\xi$。
\end{exercise}
\begin{exercise}{568}{常规·积分平均值}
设 $f$ 在 $[0,1]$ 上连续且单调减少。证明 $\int_0^xf(t)\dd t\ge xf(x)$（$0<x\le1$）。提示：比较 $0\le t\le x$ 时的 $f(t)$ 与 $f(x)$。
\end{exercise}
\begin{exercise}{569}{基础·函数平均值}
求 $f(x)=\sqrt{4-x^2}$ 在 $[-2,2]$ 上的平均值。
\end{exercise}
\begin{exercise}{570}{基础·黎曼和变式}
把 $\lim_{n\to\infty}\sum_{k=1}^n\frac{k}{n^2+k^2}$ 写成定积分，并用换元法求值。
\end{exercise}
\begin{exercise}{571}{常规·积分不等式}
设 $f$ 在 $[0,1]$ 上连续且 $0\le f\le1$。证明 $0\le\int_0^1f^2(x)\dd x\le\int_0^1f(x)\dd x$。提示：先比较被积函数。
\end{exercise}
\subsection{微积分基本公式}
\begin{exercise}{572}{基础·双变限求导}
求 $y=\int_x^{x^2}\sin(t^2)\dd t$ 的导数。
\end{exercise}
\begin{exercise}{573}{基础·积分上限求导}
设 $y=\int_0^{x^2}\sqrt{1+t^2}\dd t$，求 $y'$ 及 $y'(1)$。
\end{exercise}
\begin{exercise}{574}{基础·双变限求导}
求 $y=\int_x^{2x}\frac{\dd t}{1+t^2}$ 的导数及 $y'(0)$。
\end{exercise}
\begin{exercise}{575}{基础·变限积分极限}
求 $\lim_{x\to0}\frac{\int_0^x\cos(t^2)\dd t}{x}$。
\end{exercise}
\begin{exercise}{576}{基础·变限积分极限}
利用洛必达法则求 $\lim_{x\to0}\frac{\int_0^x\arctan t\dd t}{x^2}$。
\end{exercise}
\begin{exercise}{577}{常规·积分商极限}
求 $\lim_{x\to0}\frac{(\int_0^xe^{t^2}\dd t)^2}{x^2}$。提示：先求积分与 $x$ 的比值极限。
\end{exercise}
\begin{exercise}{578}{常规·含参变限积分}
设 $f$ 在原点附近连续，求 $\lim_{x\to0}\frac{\int_0^x tf(t)\dd t}{x^2}$，结果用 $f(0)$ 表示。
\end{exercise}
\begin{exercise}{579}{基础·无穷远积分极限}
利用洛必达法则求 $\lim_{x\to+\infty}\frac{\int_0^x\arctan t\dd t}{x}$。
\end{exercise}
\begin{exercise}{580}{常规·隐式变限积分}
设 $\int_0^ye^t\dd t+\int_0^x\cos t\dd t=0$ 在原点附近确定 $y=y(x)$。求 $y'(0)$。
\end{exercise}
\begin{exercise}{581}{基础·变限函数单调性}
设 $f(x)=\int_0^xe^{-\frac{t^2}{2}}\dd t$，求 $f'(x)$ 并判断 $f$ 的单调性。
\end{exercise}
\begin{exercise}{582}{常规·求导恢复数据}
设 $f$ 连续，$\varphi(x)=x\int_0^{x^2}f(t)\dd t$，已知 $\varphi(1)=1$、$\varphi'(1)=5$。求 $f(1)$。
\end{exercise}
\begin{exercise}{583}{常规·积分函数零点}
设 $F(x)=\int_0^x(1+t^2)\dd t-1$。证明 $F$ 在 $(0,1)$ 内恰有一个零点。提示：计算端点值和导数。
\end{exercise}
\begin{exercise}{584}{基础·牛顿莱布尼茨公式}
计算 $\int_{-1}^{\frac{1}{2}}e^{2x}\dd x$。
\end{exercise}
\begin{exercise}{585}{基础·基本公式}
计算 $\int_0^2\frac{x}{\sqrt{1+x^2}}\dd x$。
\end{exercise}
\begin{exercise}{586}{基础·部分分式积分}
计算 $\int_0^1\frac{\dd x}{x^2+5x+6}$。
\end{exercise}
\begin{exercise}{587}{基础·绝对值积分}
计算 $\int_{-1}^3|2-x|\dd x$，先确定绝对值内表达式改变符号的位置。
\end{exercise}
\begin{exercise}{588}{基础·三角绝对值积分}
计算 $\int_0^{2\pi}|\sin x|\dd x$。
\end{exercise}
\begin{exercise}{589}{常规·分段变限积分}
设 $f(x)=2x$（$0\le x\le1$），$f(x)=3$（$1<x\le2$）。求 $P(x)=\int_0^xf(t)\dd t$ 的分段表达式。
\end{exercise}
\begin{exercise}{590}{常规·积分常数恢复}
连续函数 $f$ 在 $[0,1]$ 上满足 $f(x)=x+\frac12\int_0^1f(t)\dd t$。令 $A=\int_0^1f(t)\dd t$，求 $A$ 和 $f(x)$。
\end{exercise}
\begin{exercise}{591}{常规·积分方程}
连续函数 $f$ 在 $[0,1]$ 上满足 $f(x)=x\int_0^1f(t)\dd t-1$。求 $\int_0^1f(t)\dd t$ 和 $f(x)$。
\end{exercise}
\begin{exercise}{592}{基础·复合换元}
计算 $\int_1^2(2x+1)(x^2+x)\dd x$，提示：令 $u=x^2+x$。
\end{exercise}
\begin{exercise}{593}{基础·变限函数极值}
设 $F(x)=\int_0^x(t-1)e^{-t}\dd t$，$x\ge0$。求 $F$ 的单调区间及极小值。
\end{exercise}
\begin{exercise}{594}{常规·变限积分小综合}
设 $F(x)=\int_0^{x^2}(1+t)\dd t$。求 $F'(x)$，并计算 $\int_0^1F'(x)\dd x$。
\end{exercise}
\subsection{定积分的换元积分法和分部积分法}
\begin{exercise}{595}{基础·根式换元}
计算 $\int_4^9\frac{\dd x}{\sqrt x-1}$，提示：令 $u=\sqrt x$。
\end{exercise}
\begin{exercise}{596}{基础·根式换元}
计算 $\int_0^1x\sqrt{3-2x}\dd x$。
\end{exercise}
\begin{exercise}{597}{基础·对数换元}
计算 $\int_1^e\frac{\dd x}{x\sqrt{1+\ln x}}$。
\end{exercise}
\begin{exercise}{598}{基础·指数换元}
计算 $\int_0^1\frac{e^x}{\sqrt{1+e^x}}\dd x$。
\end{exercise}
\begin{exercise}{599}{基础·三角换元}
计算 $\int_0^{\frac{\sqrt2}{3}}\sqrt{2-9x^2}\dd x$。
\end{exercise}
\begin{exercise}{600}{基础·三角换元}
计算 $\int_0^1\frac{\dd x}{\sqrt{1-x^2}}$，按上端点极限理解积分。
\end{exercise}
\begin{exercise}{601}{基础·根式换元}
计算 $\int_1^2x\sqrt{x^2-1}\dd x$。
\end{exercise}
\begin{exercise}{602}{常规·根式与对称性}
利用偶函数性质及三角换元，计算 $\int_{-1}^1\frac{\dd x}{1+\sqrt{1-x^2}}$。
\end{exercise}
\begin{exercise}{603}{基础·配方与几何意义}
将 $2x-x^2$ 配方，计算 $\int_0^1\sqrt{2x-x^2}\dd x$。
\end{exercise}
\begin{exercise}{604}{基础·指数换元}
计算 $\int_0^{\ln2}\frac{e^{2x}}{1+e^x}\dd x$，提示：令 $u=e^x$。
\end{exercise}
\begin{exercise}{605}{基础·平移分段积分}
设 $f(t)=1+t^2$（$t<0$），$f(t)=e^{-t}$（$t\ge0$）。计算 $\int_1^3f(x-2)\dd x$。
\end{exercise}
\begin{exercise}{606}{基础·平移与奇偶性}
设 $f(t)=te^{t^2}$（$-1\le t\le1$）。求 $\int_0^2f(x-1)\dd x$。
\end{exercise}
\begin{exercise}{607}{基础·积分函数恢复}
设 $f$ 连续且 $\int_0^xf(t)\dd t=e^{-2x}-1$。求 $f(x)$ 及 $\int_0^1f(x)\dd x$。
\end{exercise}
\begin{exercise}{608}{常规·二次积分求导}
设 $f$ 连续且 $\int_0^x(x-t)f(t)\dd t=e^x-x-1$。对等式求导两次，求 $f(x)$。
\end{exercise}
\begin{exercise}{609}{基础·奇函数积分}
计算 $\int_{-\frac{\pi}{2}}^{\frac{\pi}{2}}\frac{x}{1+\cos x}\dd x$。
\end{exercise}
\begin{exercise}{610}{基础·对数奇偶性}
利用 $(\sqrt{1+x^2}+x)(\sqrt{1+x^2}-x)=1$，计算 $\int_{-1}^1\ln(x+\sqrt{1+x^2})\dd x$。
\end{exercise}
\begin{exercise}{611}{常规·对称换元}
用 $x\mapsto-x$ 计算 $\int_{-1}^1\frac{\dd x}{1+e^x}$。
\end{exercise}
\begin{exercise}{612}{常规·反三角对称性}
利用 $\arccos(-x)=\pi-\arccos x$，计算 $\int_{-1}^1\cos x\arccos x\dd x$。
\end{exercise}
\begin{exercise}{613}{基础·奇偶分解}
计算 $\int_{-1}^1(x^4+\sin x)\dd x$。
\end{exercise}
\begin{exercise}{614}{基础·周期与绝对值}
计算 $\int_0^{2\pi}\sqrt{1-\sin^2x}\dd x$，注意根式等于 $|\cos x|$。
\end{exercise}
\begin{exercise}{615}{基础·三角积分}
计算 $\int_0^{\frac{\pi}{4}}\sin x\cos x\dd x$。
\end{exercise}
\begin{exercise}{616}{常规·对称换元}
利用 $x\mapsto\frac{\pi}{2}-x$，计算 $\int_0^{\frac{\pi}{2}}\frac{\sin x}{\sin x+\cos x}\dd x$。
\end{exercise}
\begin{exercise}{617}{基础·分部积分}
计算 $\int_0^1\ln(1+x^2)\dd x$。
\end{exercise}
\begin{exercise}{618}{基础·反三角分部积分}
计算 $\int_0^{\frac{1}{2}}\arcsin x\dd x$。
\end{exercise}
\begin{exercise}{619}{常规·反三角分部积分}
用分部积分法计算 $\int_{-1}^1x\arccos x\dd x$。
\end{exercise}
\begin{exercise}{620}{基础·对数分部积分}
计算 $\int_0^1x\ln(1+x)\dd x$。
\end{exercise}
\begin{exercise}{621}{常规·循环分部积分}
计算 $\int_0^{\frac{\pi}{2}}e^{-x}\sin2x\dd x$。
\end{exercise}
\begin{exercise}{622}{常规·换元与分部积分}
计算 $\int_0^1e^{-\sqrt x}\dd x$，先令 $t=\sqrt x$。
\end{exercise}
\begin{exercise}{623}{基础·三角换元}
计算 $\int_0^1\frac{\arcsin x}{\sqrt{1-x^2}}\dd x$，按上端点极限理解积分。
\end{exercise}
\begin{exercise}{624}{常规·反三角分部积分}
计算 $\int_0^1\arctan x\dd x$。
\end{exercise}
\begin{exercise}{625}{基础·倒数换元}
计算 $\int_1^2x^{-2}e^{\frac{1}{x}}\dd x$。
\end{exercise}
\begin{exercise}{626}{常规·根式换元与分部积分}
计算 $\int_0^{\pi^2}\cos\sqrt x\dd x$，先令 $t=\sqrt x$。
\end{exercise}
\begin{exercise}{627}{基础·端点数据积分}
设 $f\in C^2[0,2]$，$f(0)=1$、$f(2)=3$、$f'(2)=5$。求 $\int_0^2xf''(x)\dd x$。
\end{exercise}
\begin{exercise}{628}{基础·周期变限积分}
设 $F(x)=\int_x^{x+2\pi}\sin t\dd t$。求 $F'(x)$ 及 $F(x)$。
\end{exercise}
\begin{exercise}{629}{常规·已知原函数}
已知 $F(x)=e^{x^2}$ 是 $f(x)$ 的一个原函数。计算 $\int_0^1xf'(x)\dd x$，提示：先分部积分。
\end{exercise}
\begin{exercise}{630}{常规·积分常数恢复}
设连续函数 $f$ 满足 $f(x)=x+\int_0^1tf(t)\dd t$。求 $f(x)$。提示：把积分记为常数后代回求值。
\end{exercise}
\begin{exercise}{631}{常规·变限积分分部计算}
设 $f(x)=\int_0^x\sin t\dd t$。计算 $\int_0^\pi f(x)\dd x$。
\end{exercise}
\begin{exercise}{632}{常规·变限函数分部积分}
设 $f(x)=\int_0^xe^{-t^2}\dd t$。计算 $\int_0^1xf(x)\dd x$，结果允许含 $\int_0^1e^{-t^2}\dd t$。
\end{exercise}
\begin{exercise}{633}{常规·求导与基本公式}
设 $f(x)=\int_0^x(1+t^2)\dd t$。计算 $\int_0^1xf'(x)\dd x$，并用分部积分核验。
\end{exercise}
\begin{exercise}{634}{基础·对称换元变式}
用 $x\mapsto\frac{\pi}{2}-x$ 计算 $\int_0^{\frac{\pi}{2}}\frac{\cos x}{\sin x+\cos x}\dd x$。
\end{exercise}
\begin{exercise}{635}{常规·分部积分小综合}
设 $I_n=\int_0^1x^ne^x\dd x$，$n\ge1$ 为整数。用分部积分得到 $I_n=e-nI_{n-1}$，并据此求 $I_2$。
\end{exercise}
\subsection{反常积分}
\begin{exercise}{636}{基础·无穷积分判敛}
分别判断 $\int_1^{+\infty}x^{-\frac{1}{2}}\dd x$ 与 $\int_1^{+\infty}x^{-2}\dd x$ 的敛散性，对收敛者求值。
\end{exercise}
\begin{exercise}{637}{基础·对数反常积分}
计算 $\int_2^{+\infty}\frac{\dd x}{x\ln^2x}$。
\end{exercise}
\begin{exercise}{638}{基础·全轴积分}
计算 $\int_{-\infty}^{+\infty}\frac{\dd x}{x^2+4x+5}$，先对分母配方。
\end{exercise}
\begin{exercise}{639}{基础·部分分式反常积分}
计算 $\int_1^{+\infty}\frac{\dd x}{x^2+x}$，先在 $[1,R]$ 上积分再取极限。
\end{exercise}
\begin{exercise}{640}{基础·根式反常积分}
计算 $\int_1^{+\infty}x^{-\frac{3}{2}}\dd x$。
\end{exercise}
\begin{exercise}{641}{常规·全轴换元积分}
令 $u=e^x$，计算 $\int_{-\infty}^{+\infty}\frac{e^x}{1+e^{2x}}\dd x$。
\end{exercise}
\begin{exercise}{642}{基础·指数反常积分}
计算 $\int_0^{+\infty}e^{-2x}\dd x$。
\end{exercise}
\begin{exercise}{643}{常规·反常分部积分}
计算 $\int_1^{+\infty}\frac{\ln x}{x^2}\dd x$，写出分部积分中无穷端边界项的极限。
\end{exercise}
\begin{exercise}{644}{常规·三角代换}
计算 $\int_0^{+\infty}\frac{\dd x}{(1+x^2)^2}$，提示：令 $x=\tan t$。
\end{exercise}
\begin{exercise}{645}{基础·反三角反常积分}
计算 $\int_1^{+\infty}\frac{\dd x}{1+x^2}$。
\end{exercise}
\begin{exercise}{646}{常规·反常分部积分}
计算 $\int_0^{+\infty}xe^{-x}\dd x$。
\end{exercise}
\begin{exercise}{647}{基础·全轴绝对收敛积分}
判断 $\int_{-\infty}^{+\infty}xe^{-x^2}\dd x$ 是否绝对收敛，并求其值。提示：分别考察两个半轴。
\end{exercise}
\begin{exercise}{648}{基础·比较判别法}
利用 $\ln x<x$（$x\ge2$），判断 $\int_2^{+\infty}\frac{\dd x}{\ln x}$ 的敛散性。
\end{exercise}
\begin{exercise}{649}{常规·高斯型积分判敛}
用 $x\ge1$ 时 $e^{-x^2}\le e^{-x}$，证明 $\int_0^{+\infty}e^{-x^2}\dd x$ 收敛，无需求准确值。
\end{exercise}
\begin{exercise}{650}{基础·指数反常积分}
计算 $\int_{-\infty}^0e^{2t}\dd t$。
\end{exercise}
\begin{exercise}{651}{基础·瑕积分判敛}
判断 $\int_0^1x^{-\frac{1}{2}}\dd x$ 与 $\int_0^1x^{-2}\dd x$ 的敛散性，对收敛者求值。
\end{exercise}
\begin{exercise}{652}{常规·对数瑕积分}
计算 $\int_0^1x\ln x\dd x$，写出 $x\to0^+$ 时边界项的极限。
\end{exercise}
\begin{exercise}{653}{基础·根式瑕积分}
计算 $\int_0^2\frac{\dd x}{\sqrt{4-x^2}}$，用极限处理上端点。
\end{exercise}
\begin{exercise}{654}{基础·瑕点识别}
判断 $\int_0^1\frac{\dd x}{x(\ln x)^2}$ 的敛散性，分别考察 $0$ 和 $1$ 两个端点。
\end{exercise}
\begin{exercise}{655}{基础·幂次判别}
讨论 $\int_0^1x^{-p}\dd x$ 随实参数 $p$ 的敛散性，收敛时求值。
\end{exercise}
\begin{exercise}{656}{常规·倒数换元}
计算 $\int_1^{+\infty}\frac{\ln x}{x^2}\dd x$，这次使用 $t=\frac{1}{x}$ 将其化为瑕积分。
\end{exercise}
\begin{exercise}{657}{基础·分段反常积分}
设 $f(x)=x^{-\frac{1}{2}}$（$0<x\le1$），$f(x)=x^{-2}$（$x>1$）。计算 $\int_0^{+\infty}f(x)\dd x$。
\end{exercise}
\begin{exercise}{658}{常规·内部瑕点}
判断 $\int_{-1}^1x^{-2}\dd x$ 是否存在，先在 $0$ 处分成两个反常积分。
\end{exercise}
\begin{exercise}{659}{基础·双端根式瑕积分}
令 $x=\sin^2t$，计算 $\int_0^1\frac{\dd x}{\sqrt{x(1-x)}}$。
\end{exercise}
\begin{exercise}{660}{常规·对数换元小综合}
判断 $\int_e^{+\infty}\frac{\dd x}{x\ln x}$ 与 $\int_e^{+\infty}\frac{\dd x}{x(\ln x)^2}$ 的敛散性，对收敛者求值。
\end{exercise}

\clearpage
\section{第6章 定积分的应用}
\subsection{定积分在几何学上的应用}
\begin{exercise}{661}{基础·曲边梯形面积}
求 $y=\sqrt x$、$y=0$ 与 $x=4$ 围成区域的面积。
\end{exercise}
\begin{exercise}{662}{基础·横向切片面积}
求由 $x=y^2$、$x=0$、$y=1$ 围成区域的面积，使用关于 $y$ 的积分。
\end{exercise}
\begin{exercise}{663}{常规·面积与旋转体积}
求 $y=x^2$ 与 $y=x$ 围成区域的面积，并求该区域绕 $x$ 轴旋转一周的体积。
\end{exercise}
\begin{exercise}{664}{基础·抛物线间面积}
求 $y=4-x^2$ 与 $y=4-2x$ 围成区域的面积，先求交点。
\end{exercise}
\begin{exercise}{665}{常规·面积反求参数}
曲线 $y=\sqrt{ax}$（$a>0$）、$x$ 轴及 $x=3$ 围成区域的面积为 $6$。求 $a$。
\end{exercise}
\begin{exercise}{666}{基础·指数曲线面积}
求 $y=e^x$、$y=0$、$x=0$、$x=1$ 围成区域的面积。
\end{exercise}
\begin{exercise}{667}{基础·有向积分与面积}
求 $y=x^2-x$ 与 $x$ 轴围成区域的面积。
\end{exercise}
\begin{exercise}{668}{常规·切线与面积}
求 $y=x^2$ 在 $(1,1)$ 处的切线，并求切线、抛物线与 $y$ 轴围成区域的面积。
\end{exercise}
\begin{exercise}{669}{基础·横向切片}
求 $x=4-y^2$ 与 $x=0$ 围成区域的面积。
\end{exercise}
\begin{exercise}{670}{基础·三角曲线面积}
求 $y=\sin x$（$0\le x\le\pi$）与 $x$ 轴围成区域的面积。
\end{exercise}
\begin{exercise}{671}{常规·对称区域面积}
求 $y=e^{-|x|}$、$y=0$、$x=-1$、$x=1$ 围成区域的面积。
\end{exercise}
\begin{exercise}{672}{基础·对数曲线面积}
求 $y=\ln x$、$y=0$ 与 $x=e$ 围成区域的面积。
\end{exercise}
\begin{exercise}{673}{常规·两三角曲线间面积}
求 $y=\sin x$、$y=\cos x$、$x=0$、$x=\frac{\pi}{2}$ 围成区域的面积。
\end{exercise}
\begin{exercise}{674}{基础·圆弓形面积}
求单位圆 $x^2+y^2\le1$ 位于 $x$ 轴上方的区域面积，用定积分表示并计算。
\end{exercise}
\begin{exercise}{675}{基础·水平直线截面积}
求 $y=4-x^2$ 与 $y=3$ 围成区域的面积。
\end{exercise}
\begin{exercise}{676}{基础·椭圆面积}
用换元法求椭圆 $\frac{x^2}{4}+y^2=1$ 内部区域的面积。
\end{exercise}
\begin{exercise}{677}{基础·抛物线面积}
求 $y=2x-x^2$ 与 $x$ 轴围成区域的面积。
\end{exercise}
\begin{exercise}{678}{常规·分段面积}
求 $y=x^2$ 与 $y=\frac{1}{4}$ 在 $0\le x\le1$ 之间所夹区域的总面积。
\end{exercise}
\begin{exercise}{679}{基础·两抛物线面积}
求 $y=x^2$ 与 $y=2-x^2$ 围成区域的面积。
\end{exercise}
\begin{exercise}{680}{常规·面积随参数变化}
设 $a>0$，求 $y=ax$ 与 $y=x^2$ 围成区域的面积；面积为 $\frac{1}{6}$ 时求 $a$。
\end{exercise}
\begin{exercise}{681}{基础·极坐标面积}
用极坐标面积公式求曲线 $r=2\cos\theta$ 围成圆盘的面积，其中 $-\frac{\pi}{2}\le\theta\le\frac{\pi}{2}$。
\end{exercise}
\begin{exercise}{682}{常规·极坐标面积}
求心形线 $r=1+\cos\theta$（$0\le\theta\le2\pi$）围成区域的面积。
\end{exercise}
\begin{exercise}{683}{常规·参数曲线面积}
摆线一拱为 $x=t-\sin t$、$y=1-\cos t$，$0\le t\le2\pi$。求它与 $x$ 轴围成区域的面积。
\end{exercise}
\begin{exercise}{684}{基础·参数曲线面积}
曲线由 $x=t^2$、$y=t$（$0\le t\le1$）给出。求该曲线与 $x$ 轴、$x=1$ 围成区域的面积。
\end{exercise}
\begin{exercise}{685}{基础·圆盘法体积}
求 $y=\sqrt x$、$x=0$、$x=4$、$y=0$ 围成区域绕 $x$ 轴旋转一周的体积。
\end{exercise}
\begin{exercise}{686}{常规·圆环法体积}
求 $y=x$ 与 $y=x^2$ 围成区域绕 $x$ 轴旋转一周的体积，注明内、外半径。
\end{exercise}
\begin{exercise}{687}{基础·无界旋转体积}
求 $D=\{(x,y):x\ge0,0\le y\le e^{-x}\}$ 绕 $x$ 轴旋转所得体积，用反常积分表示。
\end{exercise}
\begin{exercise}{688}{基础·柱壳法体积}
求 $y=e^{-x^2}$、$x=-1$、$x=0$ 与 $x$ 轴围成区域绕 $y$ 轴旋转的体积，注意柱壳半径为 $-x$。
\end{exercise}
\begin{exercise}{689}{基础·圆盘旋转}
求圆盘 $x^2+y^2\le4$ 绕 $x$ 轴旋转所得球体的体积，先写出垂直于 $x$ 轴的截面面积。
\end{exercise}
\begin{exercise}{690}{基础·对数区域旋转}
求 $y=\ln x$、$y=0$ 与 $x=e$ 围成区域绕 $y$ 轴旋转的体积。
\end{exercise}
\begin{exercise}{691}{常规·反常积分体积}
求 $y=xe^{-x}$（$x\ge0$）与 $x$ 轴间区域绕 $x$ 轴旋转所得体积。
\end{exercise}
\begin{exercise}{692}{基础·半圆旋转体积}
求 $y=\sqrt{1-x^2}$（$-1\le x\le1$）与 $x$ 轴围成区域绕 $x$ 轴旋转的体积。
\end{exercise}
\begin{exercise}{693}{常规·平移旋转轴}
求矩形 $0\le x\le2$、$0\le y\le1$ 绕直线 $y=2$ 旋转一周的体积，注明内、外半径。
\end{exercise}
\begin{exercise}{694}{常规·切线区域旋转}
曲线 $y=x^2$ 在 $(1,1)$ 处的切线为 $y=2x-1$。求由该切线、曲线与 $y$ 轴围成区域绕 $y$ 轴旋转一周的体积，提示：使用柱壳法，$0\le x\le1$。
\end{exercise}
\begin{exercise}{695}{常规·累积体积}
设 $V(a)$ 为 $y=e^{-x}$、$y=0$、$x=0$、$x=a>0$ 围成区域绕 $x$ 轴旋转的体积。求 $V(a)$，并求满足 $V(a)=\frac{\pi}{4}$ 的 $a$。
\end{exercise}
\begin{exercise}{696}{常规·旋转曲面面积}
求直线段 $y=x$（$0\le x\le1$）绕 $x$ 轴旋转所得曲面的侧面积。
\end{exercise}
\begin{exercise}{697}{常规·旋转曲面面积}
求曲线 $y=\sqrt{2x+1}$（$0\le x\le1$）绕 $x$ 轴旋转所得曲面的面积。
\end{exercise}
\begin{exercise}{698}{基础·直角坐标弧长}
求曲线 $y=x^{\frac{3}{2}}$（$0\le x\le4$）的弧长。
\end{exercise}
\begin{exercise}{699}{基础·三角对数弧长}
求曲线 $y=\ln\cos x$（$0\le x\le\frac{\pi}{3}$）的弧长。
\end{exercise}
\begin{exercise}{700}{常规·变限函数弧长}
求曲线 $y=\int_0^x\tan t\dd t$（$0\le x\le\frac{\pi}{4}$）的弧长。
\end{exercise}
\begin{exercise}{701}{基础·参数曲线弧长}
求 $x=2\cos t$、$y=2\sin t$（$0\le t\le\frac{\pi}{2}$）所表示圆弧的长度。
\end{exercise}
\begin{exercise}{702}{基础·极坐标弧长}
求对数螺线 $r=e^\theta$（$0\le\theta\le\pi$）的弧长。
\end{exercise}
\begin{exercise}{703}{常规·参数弧长}
求曲线 $x=t^2$、$y=\frac23t^3$（$0\le t\le1$）的弧长。
\end{exercise}
\begin{exercise}{704}{常规·封闭区域周长}
求 $y=x$、$y=0$ 与 $x=1$ 围成三角形的周长，并用弧长积分核验斜边长度。
\end{exercise}
\begin{exercise}{705}{基础·曲率半径}
利用 $\rho=\frac{(1+(y')^2)^{\frac{3}{2}}}{|y''|}$，求抛物线 $y=x^2$ 在原点的曲率半径。
\end{exercise}
\begin{exercise}{706}{基础·轴外柱壳体积}
求矩形 $0\le x\le1$、$0\le y\le2$ 绕直线 $x=2$ 旋转一周的体积。
\end{exercise}
\begin{exercise}{707}{基础·绝对值区域旋转}
求 $y=1-|x|$ 与 $x$ 轴围成区域绕 $x$ 轴旋转所得体积。
\end{exercise}
\begin{exercise}{708}{常规·已知截面求体积}
某立体在 $0\le x\le2$ 内垂直于 $x$ 轴的截面均为正方形，边长为 $\sqrt{4-x^2}$。求立体体积。
\end{exercise}
\begin{exercise}{709}{基础·平移轴旋转变式}
求 $y=x$、$y=0$ 与 $x=1$ 围成区域绕直线 $y=2$ 旋转的体积。
\end{exercise}
\begin{exercise}{710}{基础·旋转体积变式}
求 $y=x^2$ 与 $y=2x$ 围成区域绕 $x$ 轴旋转的体积。
\end{exercise}
\begin{exercise}{711}{常规·有限旋转体积}
求 $D=\{(x,y):x\ge1,0\le y\le\frac{1}{x}\}$ 绕 $x$ 轴旋转所得体积，先截断到 $x=R$ 再取极限。
\end{exercise}
\subsection{定积分在物理学上的应用}
\begin{exercise}{712}{基础·速度与路程}
质点在 $0\le t\le2$ 内的速度为 $v(t)=3t^2+1$ 米/秒。求这段时间的路程。
\end{exercise}
\begin{exercise}{713}{常规·缆绳做功}
长 $10$ 米的均匀缆绳竖直悬挂，每米重 $20$ 牛。把缆绳全部缓慢拉到上端，求克服重力所做的功。
\end{exercise}
\begin{exercise}{714}{基础·矩形闸门压力}
竖直矩形闸门宽 $2$ 米、高 $3$ 米，上边位于水面。水密度为 $\rho$、重力加速度为 $g$。用水平薄条积分求单侧静水总压力。
\end{exercise}
\begin{exercise}{715}{基础·变力做功}
弹簧伸长 $x$ 米时弹力大小为 $100x$ 牛。将弹簧从自然长度拉长到 $0.2$ 米，求克服弹力所做的功。
\end{exercise}
\begin{exercise}{716}{基础·变密度细棒质量}
长 $2$ 米的细棒在位置 $x$ 米处的线密度为 $1+x$ 千克/米（$0\le x\le2$）。求细棒质量。
\end{exercise}
\begin{exercise}{717}{常规·圆柱容器抽水功}
半径 $1$ 米、高 $2$ 米的竖直圆柱水箱盛满水。求将水全部抽至箱顶所做的功，水密度为 $1000$ 千克/立方米，重力加速度为 $g$。
\end{exercise}
\begin{exercise}{718}{常规·均匀薄片形心}
均匀薄片为 $0\le x\le1$、$0\le y\le x$。求其面积及形心横坐标，提示：$\bar x=\frac{\int_0^1x\cdot x\dd x}{\int_0^1x\dd x}$。
\end{exercise}
\begin{exercise}{719}{基础·变力做功变式}
物体沿 $x$ 轴正向移动，所受力为 $F(x)=2x+3$ 牛。求物体从 $x=0$ 米移到 $x=2$ 米时该力所做的功。
\end{exercise}
\begin{exercise}{720}{常规·三角形水压小综合}
竖直等腰三角形薄板高 $h$、底边长 $b$，顶点在水面，底边水平且位于顶点下方。水密度为 $\rho$、重力加速度为 $g$。求薄板单侧静水总压力，提示：深度 $y$ 处的宽度为 $\frac{by}{h}$。
\end{exercise}

\clearpage
\section{第7章 微分方程}
\subsection{微分方程的基本概念}
\begin{exercise}{721}{基础·微分方程辨识}
判断 $x^2+y^2=1$、$y'+y=x$、$y''+y=0$ 中哪些为微分方程，并指出其阶数。
\end{exercise}
\begin{exercise}{722}{基础·阶数与线性}
判断 $y''+(y')^2+y=0$ 的阶数及是否为线性方程。
\end{exercise}
\begin{exercise}{723}{基础·标准形式}
将 $x\dd y+(x+y)\dd x=0$ 在 $x>0$ 上写成一阶线性方程的标准形式。
\end{exercise}
\begin{exercise}{724}{基础·解的验证}
验证 $y=5x^2$ 是 $xy'-2y=0$ 的解，并求它在 $x=1$ 处的函数值。
\end{exercise}
\begin{exercise}{725}{基础·初值确定常数}
已知 $y=C_1\cos x+C_2\sin x$ 是 $y''+y=0$ 的通解，求满足 $y(0)=2$、$y'(0)=3$ 的解。
\end{exercise}
\begin{exercise}{726}{基础·通解与特解}
验证 $y=x+C$ 是 $y'=1$ 的通解，并求经过 $(0,2)$ 的特解。
\end{exercise}
\begin{exercise}{727}{基础·隐式解验证}
在 $y>0$ 的分支上，验证 $x^2+y^2=4$ 满足 $y'=-x/y$。
\end{exercise}
\begin{exercise}{728}{基础·常值解变式}
分别检验 $y\equiv0$ 与 $y\equiv1$ 是否为 $y'=y(1-y)$ 的解。
\end{exercise}
\begin{exercise}{729}{常规·消去任意常数}
消去 $y=C_1e^x+C_2e^{-x}$ 中的任意常数，建立二阶微分方程，并求满足 $y(0)=1$、$y'(0)=0$ 的解。
\end{exercise}
\subsection{可分离变量的微分方程}
\begin{exercise}{730}{基础·分离变量辨识}
分别将 $y'=xy$ 与 $y'=x/(1+y^2)$ 分离变量；对第一个方程另外检查零解。
\end{exercise}
\begin{exercise}{731}{基础·可分离变量方程}
求 $y'=y(1-x)/x$ 在 $x>0$ 上满足 $y(1)=e^{-1}$ 的解。
\end{exercise}
\begin{exercise}{732}{基础·分离变量初值}
求 $y'=\sin x\cos x/y^2$、$y(0)=1$ 的解，取包含初始点的 $y>0$ 分支。
\end{exercise}
\begin{exercise}{733}{基础·指数分离变量}
求 $(1+e^x)y'-ye^x=0$、$y(0)=2$ 的解。
\end{exercise}
\begin{exercise}{734}{基础·根式初值解}
求 $y'=x/y$、$y(0)=1$ 的正值解。
\end{exercise}
\begin{exercise}{735}{基础·隐式分离变量}
求 $e^{x-y}+yy'=0$、$y(0)=1$ 在初始点附近的隐式解。
\end{exercise}
\begin{exercise}{736}{常规·斜率恢复曲线}
曲线经过 $(0,-1)$，其在 $(x,y)$ 处的切线斜率为 $y^2/(1+x^2)$。求初始点附近的曲线方程。
\end{exercise}
\begin{exercise}{737}{常规·积分方程求导}
设连续函数 $f$ 满足 $f(x)=\int_0^xf(t)\dd t+2$。求 $f(x)$，先由等式恢复初值。
\end{exercise}
\begin{exercise}{738}{基础·可分离变量初值}
求 $y'=2xy$、$y(0)=1$ 的解。
\end{exercise}
\begin{exercise}{739}{常规·积分形式初值}
设连续函数 $f$ 满足 $f(x)=1+2\int_0^xf(t)\dd t$。求 $f(x)$ 并代回核验。
\end{exercise}
\begin{exercise}{740}{基础·常值解与分离变量}
求 $y'=-2y$ 的通解，包括零解。
\end{exercise}
\begin{exercise}{741}{常规·积分条件定常数}
设 $f'=xf/(1+x^2)$、$f(0)=a$，且 $\int_0^{\sqrt3}xf(x)\dd x=7$。求 $a$ 与 $f(x)$。
\end{exercise}
\begin{exercise}{742}{基础·平移后分离变量}
求 $y'=(y+1)^2$、$y(0)=0$ 在原点附近的解，提示：令 $u=y+1$。
\end{exercise}
\begin{exercise}{743}{基础·分离变量隐式解}
求 $(1+y^2)y'=x$ 的通解。
\end{exercise}
\begin{exercise}{744}{基础·可分离变量变式}
求 $y'=y^2$、$y(0)=1$ 在原点附近的解。
\end{exercise}
\begin{exercise}{745}{常规·增长模型}
种群数量满足 $N'=N(1-N/100)$、$N(0)=20$，$t\ge0$。求 $N(t)$，并计算 $N(t)=50$ 时的 $t$。
\end{exercise}
\subsection{齐次方程}
\begin{exercise}{746}{基础·齐次方程代换}
将 $y'=1+y/x$（$x>0$）用 $u=y/x$ 化为可分离变量方程，并求通解。
\end{exercise}
\begin{exercise}{747}{基础·根式齐次方程}
求 $y'=2\sqrt{y/x}+y/x$ 满足 $y(1)=1$ 的解，只考虑初始点附近的 $x>0,y>0$ 分支。
\end{exercise}
\begin{exercise}{748}{基础·齐次代换}
求 $y'=y/x+(y/x)^2$ 满足 $y(1)=1$ 的解，取初始点附近 $x>0$ 的分支。
\end{exercise}
\begin{exercise}{749}{常规·齐次初值问题}
求 $y'=(x^2+y^2)/(xy)$、$y(1)=2$ 的正值解，其中 $x>0$，只需给出初始点附近的表达式。
\end{exercise}
\begin{exercise}{750}{基础·切线斜率模型}
曲线经过 $(1,0)$，切线斜率为 $1+y/x$。求曲线在 $x>0$ 上的方程。
\end{exercise}
\begin{exercise}{751}{基础·常比例解}
检验 $y=x$ 是否为 $y'=2y/x-(y/x)^2$（$x>0$）的解，并写出其满足的初值 $y(1)$。
\end{exercise}
\begin{exercise}{752}{常规·对数齐次方程}
求 $xy'+y(\ln x-\ln y)=0$、$y(1)=e$ 的正值解，$x>0$。提示：令 $u=y/x$。
\end{exercise}
\begin{exercise}{753}{基础·齐次微分形式}
求 $(x-y)\dd x+x\dd y=0$ 在 $x>0$ 上的通解。
\end{exercise}
\begin{exercise}{754}{常规·齐次方程隐式解}
求 $y'=y/x+x/y$、$y(1)=1$ 的正值解，$x>0$。
\end{exercise}
\begin{exercise}{755}{基础·指数齐次变式}
求 $xy'=y+x$（$x>0$）满足 $y(1)=0$ 的解。
\end{exercise}
\begin{exercise}{756}{常规·齐次求解与核验}
求 $xy'=y+y^2/x$（$x>0$）、$y(1)=-1$ 的解，并代入原方程核验。
\end{exercise}
\subsection{一阶线性微分方程}
\begin{exercise}{757}{基础·齐次线性方程}
求 $y'+2xy=0$ 的通解及满足 $y(0)=1$ 的解。
\end{exercise}
\begin{exercise}{758}{基础·三角系数初值}
求 $y'=y\tan x$、$y(0)=1$ 在 $(-\pi/2,\pi/2)$ 上的解。
\end{exercise}
\begin{exercise}{759}{基础·线性方程辨识}
判断 $y'+y^2=x$、$y'+xy=x^2$、$(y')^2+y=0$ 中哪些是一阶线性方程，并对符合者写出积分因子。
\end{exercise}
\begin{exercise}{760}{基础·一阶线性初值}
求 $y'+2xy=4x$、$y(0)=1$ 的解。
\end{exercise}
\begin{exercise}{761}{基础·积分因子}
求 $y'+y\cos x=e^{-\sin x}$ 的通解。
\end{exercise}
\begin{exercise}{762}{常规·非恒定系数线性方程}
求 $(x+1)y'-y=e^x(x+1)^2$、$y(0)=2$ 在 $x>-1$ 上的解。
\end{exercise}
\begin{exercise}{763}{基础·对数外力}
求 $y'+y/x=\ln x$、$y(1)=0$ 在 $x>0$ 上的解。
\end{exercise}
\begin{exercise}{764}{常规·交换自因变量}
把 $x$ 看作 $y$ 的函数，求 $\frac{\dd x}{\dd y}+x=y$ 的通解，并求经过 $(x,y)=(0,0)$ 的隐式特解。
\end{exercise}
\begin{exercise}{765}{常规·逆函数线性方程}
求 $y\dd x+(2x-y)\dd y=0$ 在 $y>0$ 上的隐式通解，提示：以 $x$ 为 $y$ 的函数。
\end{exercise}
\begin{exercise}{766}{基础·隐式线性初值}
求 $y\dd x+(x-3y^2)\dd y=0$、$y(1)=1$ 在初始点附近的隐式解。
\end{exercise}
\begin{exercise}{767}{常规·逆函数求解}
把 $x$ 看作 $y$ 的函数，求 $y'=1/(x+y)$ 经过 $(0,1)$ 的隐式解。
\end{exercise}
\begin{exercise}{768}{基础·全微分识别}
求 $y\dd x+(x+y)\dd y=0$ 的隐式通解，提示：识别 $\dd(xy+y^2/2)$。
\end{exercise}
\begin{exercise}{769}{常规·切线斜率与线性方程}
曲线 $y=f(x)$ 经过 $(0,1)$，切线斜率为 $\sec x-f(x)\tan x$。求它在 $(-\pi/2,\pi/2)$ 上的方程。
\end{exercise}
\begin{exercise}{770}{常规·积分方程求解}
设连续函数 $f$ 满足 $f(x)=\int_0^xf(t)\dd t-x^2$。求 $f(x)$，先写出微分方程与初值。
\end{exercise}
\begin{exercise}{771}{基础·线性初值方程}
求 $y'+y=e^{-x}$、$y(0)=0$ 的解。
\end{exercise}
\begin{exercise}{772}{基础·一阶线性变式}
求 $y'+2y=1+e^{-x}$ 的通解。
\end{exercise}
\begin{exercise}{773}{常规·分段外力小综合}
系统满足 $y'+y=u(x)$、$y(0)=0$，其中 $u(x)=1$（$0\le x<1$）、$u(x)=0$（$x>1$）。分段求解，并令 $y$ 在 $x=1$ 连续。
\end{exercise}
\subsection{可降阶的高阶微分方程}
\begin{exercise}{774}{基础·二次积分}
求 $y''=x$ 的通解及满足 $y(0)=0$、$y'(0)=1$ 的解。
\end{exercise}
\begin{exercise}{775}{基础·三角指数积分}
求 $y''=e^{3x}+\sin2x$ 的通解。
\end{exercise}
\begin{exercise}{776}{基础·二阶初值方程}
求 $y''=xe^x$、$y(0)=y'(0)=0$ 的解。
\end{exercise}
\begin{exercise}{777}{基础·反三角函数积分}
求 $y''=1/(1+x^2)$、$y(0)=y'(0)=0$ 的解。
\end{exercise}
\begin{exercise}{778}{基础·指数积分初值}
求 $y''=e^{2x}$、$y(0)=y'(0)=0$ 的解。
\end{exercise}
\begin{exercise}{779}{常规·几何条件恢复初值}
曲线经过 $(0,1)$，在该点与 $y=x/2+1$ 相切，且满足 $y''=x$。求曲线方程。
\end{exercise}
\begin{exercise}{780}{基础·导数代换降阶}
求 $y''=y'+x$ 的通解，提示：令 $p=y'$。
\end{exercise}
\begin{exercise}{781}{基础·导数代换初值}
求 $y''+y'=0$、$y(0)=0$、$y'(0)=1$ 的解。
\end{exercise}
\begin{exercise}{782}{基础·变系数降阶}
求 $xy''+3y'=0$ 在 $x>0$ 上的通解，提示：令 $p=y'$。
\end{exercise}
\begin{exercise}{783}{基础·导数线性方程}
求 $y''-2y'=0$、$y(0)=0$、$y'(0)=-1$ 的解。
\end{exercise}
\begin{exercise}{784}{常规·速度方程降阶}
求 $(1+x^2)y''=2xy'$、$y(0)=1$、$y'(0)=3$ 的解，提示：先求 $p=y'$。
\end{exercise}
\begin{exercise}{785}{基础·二阶常系数降阶}
求 $y''=2y'$ 的通解。
\end{exercise}
\begin{exercise}{786}{常规·乘积求导降阶}
求 $yy''+(y')^2=0$、$y(0)=1$、$y'(0)=1/2$ 在原点附近的正值解，提示：$(yy')'=0$。
\end{exercise}
\begin{exercise}{787}{基础·简单二阶方程}
求 $y''=2$、$y(0)=1$、$y'(0)=0$ 的解。
\end{exercise}
\begin{exercise}{788}{基础·降阶线性初值}
求 $y''+2y'=2$、$y(0)=y'(0)=0$ 的解。
\end{exercise}
\begin{exercise}{789}{基础·乘积导数变式}
求 $xy''+y'=x$、$y(1)=0$、$y'(1)=1$ 在 $x>0$ 上的解，提示：$(xy')'=x$。
\end{exercise}
\begin{exercise}{790}{常规·降阶求解小综合}
求 $y''+y'=e^x$、$y(0)=0$、$y'(0)=1$ 的解，并核验两个初值。
\end{exercise}
\subsection{高阶线性微分方程}
\begin{exercise}{791}{基础·线性相关性}
判断函数组 $\{x^2,2x^2\}$ 与 $\{\sin x,\cos x\}$ 在实轴上的线性相关性。对第二组可在 $x=0,\pi/2$ 处取值。
\end{exercise}
\begin{exercise}{792}{基础·基本解组}
验证 $e^x,e^{-x}$ 是 $y''-y=0$ 的线性无关解，并写出通解。
\end{exercise}
\begin{exercise}{793}{常规·非齐次解结构}
验证 $y_p=x-1$ 是 $y'+y=x$ 的一个特解。求对应齐次方程的通解，进而写出原方程通解。
\end{exercise}
\begin{exercise}{794}{基础·冗余常数辨识}
将 $y=C_1e^x+2C_2e^x$ 化为只含一个任意常数的形式，并说明它满足哪个一阶齐次方程。
\end{exercise}
\begin{exercise}{795}{基础·朗斯基行列式}
计算 $f_1=e^x$、$f_2=e^{2x}$ 的朗斯基行列式 $W=f_1f_2'-f_2f_1'$，判断两函数是否线性无关。
\end{exercise}
\begin{exercise}{796}{基础·三角基本解}
写出 $y''+4y=0$ 的通解，并求满足 $y(0)=0$、$y'(0)=1$ 的解。
\end{exercise}
\begin{exercise}{797}{常规·通解验证}
验证 $y=C_1e^x+C_2e^{2x}+\frac1{12}e^{5x}$ 满足 $y''-3y'+2y=e^{5x}$。
\end{exercise}
\begin{exercise}{798}{常规·解的差}
设 $y_1,y_2$ 均满足 $y''+y=x$。证明 $y_1-y_2$ 满足 $y''+y=0$，提示：将两个方程相减。
\end{exercise}
\begin{exercise}{799}{基础·解的叠加变式}
验证 $y_1=\cos x$、$y_2=\sin x$ 均满足 $y''+y=0$，并求 $2y_1+3y_2$ 在 $x=0$ 处的函数值与导数值。
\end{exercise}
\begin{exercise}{800}{常规·已知解降阶}
已知 $y_1=x$ 是 $x^2y''-2xy'+2y=0$（$x>0$）的解。令 $y=xv$ 求通解，再由 $y(1)=0$、$y'(1)=1$ 确定常数。
\end{exercise}
\subsection{常系数齐次线性微分方程}
\begin{exercise}{801}{基础·常系数方程}
求 $y''+2y'=0$ 的通解。
\end{exercise}
\begin{exercise}{802}{基础·由特征根写方程}
求以 $y=C_1e^{2x}+C_2e^{-x}$ 为通解且 $y''$ 系数为 $1$ 的二阶常系数齐次方程。
\end{exercise}
\begin{exercise}{803}{基础·复特征根}
求 $y''-4y'+13y=0$ 的通解。
\end{exercise}
\begin{exercise}{804}{基础·重根恢复方程}
已知 $e^{2x}$ 与 $xe^{2x}$ 是 $y''+py'+qy=0$ 的基本解。求 $p,q$ 及通解。
\end{exercise}
\begin{exercise}{805}{基础·不同实根初值}
求 $y''+5y'+4y=0$、$y(0)=1$、$y'(0)=0$ 的解。
\end{exercise}
\begin{exercise}{806}{基础·不同实根初值}
求 $y''+4y'+3y=0$、$y(0)=2$、$y'(0)=6$ 的解。
\end{exercise}
\begin{exercise}{807}{基础·重复特征根}
求 $4y''-12y'+9y=0$ 的通解。
\end{exercise}
\begin{exercise}{808}{基础·重根初值问题}
求 $s''+2s'+s=0$、$s(0)=2$、$s'(0)=-1$ 的解，自变量为 $t$。
\end{exercise}
\begin{exercise}{809}{基础·复根通解}
求 $y''+8y'+17y=0$ 的通解。
\end{exercise}
\begin{exercise}{810}{常规·复根初值问题}
求 $y''+10y'+29y=0$、$y(0)=1$、$y'(0)=-1$ 的解。
\end{exercise}
\begin{exercise}{811}{基础·二阶特征方程}
求 $y''-y=0$ 的通解及满足 $y(0)=1$、$y'(0)=1$ 的解。
\end{exercise}
\begin{exercise}{812}{基础·二阶复根方程}
求 $y''+y=0$ 的通解及满足 $y(0)=0$、$y'(0)=2$ 的解。
\end{exercise}
\begin{exercise}{813}{常规·微分关系恢复初值}
已知 $f''+f=0$、$f(0)=1$、$f'(0)=1$。求 $f(x)$，并核验它满足 $f'(x)=f(-x)$。
\end{exercise}
\begin{exercise}{814}{基础·已知解验证}
验证 $f(x)=e^x$ 同时满足 $f''+f'-2f=0$ 与 $f''+f=2e^x$。
\end{exercise}
\begin{exercise}{815}{基础·重根方程变式}
求 $y''+4y'+4y=0$、$y(0)=0$、$y'(0)=1$ 的解。
\end{exercise}
\begin{exercise}{816}{常规·常系数初值小综合}
求 $y''+2y'+2y=0$、$y(0)=1$、$y'(0)=0$ 的解，并计算 $y(\pi/2)$。
\end{exercise}
\subsection{常系数非齐次线性微分方程}
\begin{exercise}{817}{基础·多项式特解}
求 $y''-4y'+3y=x$ 的通解，提示：可设特解 $y_p=Ax+B$。
\end{exercise}
\begin{exercise}{818}{基础·多项式外力}
求 $y''+5y'+4y=3-2x$ 的通解。
\end{exercise}
\begin{exercise}{819}{基础·零特征根}
求 $y''+y'=x$ 的通解，提示：可设特解 $y_p=Ax^2+Bx$。
\end{exercise}
\begin{exercise}{820}{常规·恢复系数与求解}
已知 $y''+ay'+by=0$ 的通解为 $(C_1+C_2x)e^x$。求 $a,b$，再求对应非齐次方程 $y''+ay'+by=x$ 的通解。
\end{exercise}
\begin{exercise}{821}{常规·单根共振}
求 $y''-5y'+6y=2e^{2x}$ 的通解，提示：设特解 $y_p=Axe^{2x}$。
\end{exercise}
\begin{exercise}{822}{基础·非共振指数外力}
求 $y''-3y'+2y=e^{3x}$ 的通解。
\end{exercise}
\begin{exercise}{823}{基础·特解系数}
对 $y''-y=3e^{2x}$，设特解 $y_p=Ae^{2x}$。求 $A$ 并写出通解。
\end{exercise}
\begin{exercise}{824}{常规·给定特解求通解}
验证 $y_p=x$ 是 $y''-y=-x$ 的特解，并求满足 $y(0)=1$、$y'(0)=0$ 的解。
\end{exercise}
\begin{exercise}{825}{常规·三角共振}
求 $y''+4y=\cos2x$ 的通解，提示：可设特解 $y_p=Ax\sin2x$。
\end{exercise}
\begin{exercise}{826}{常规·外力分解}
求 $y''+4y=\cos^2x$ 的通解，先使用 $\cos^2x=(1+\cos2x)/2$。
\end{exercise}
\begin{exercise}{827}{基础·三角非共振}
求 $y''+4y=\cos x$ 的通解。
\end{exercise}
\begin{exercise}{828}{常规·外力叠加}
求 $y''-y=e^x+4\cos x$ 的通解，分别为两项外力求特解。
\end{exercise}
\begin{exercise}{829}{基础·非共振变式}
求 $y''+y=\cos2x$ 的通解。
\end{exercise}
\begin{exercise}{830}{常规·重根共振小综合}
求 $y''-2y'+y=e^x$、$y(0)=y'(0)=0$ 的解，提示：令 $y=e^xv$。
\end{exercise}

\end{document}
